\documentclass[journal]{IEEEtran}
\usepackage[utf8]{inputenc}\usepackage[T1]{fontenc}
\usepackage{amsmath,amssymb,amsthm,booktabs,tabularx,algorithm,algpseudocode,graphicx,cite,url}
\usepackage{tikz}\usetikzlibrary{arrows.meta,positioning,fit,calc,backgrounds,decorations.pathreplacing,decorations.pathmorphing}
\newtheorem{definition}{Definition}\newtheorem{assumption}{Assumption}\newtheorem{lemma}{Lemma}\newtheorem{theorem}{Theorem}\newtheorem{corollary}{Corollary}\newtheorem{problem}{Problem}
\begin{document}
\title{Fleet-DT: A Network Digital Twin with Dynamic Bandwidth Allocation for Fleets of Unmanned Surface Vehicles}
\author{A. R. P. Domingues, C. J. D. Silva, C. A. M. Marcon, F. G. Moraes}
\maketitle
\input{abstract}
\section{Introduction}\label{sec:intro}

Unmanned surface vehicles (USVs) carry out scientific, commercial, and defense missions in oceans, coastal waters, rivers, and lakes [liu2016usv]. USVs monitor the water quality, survey the bottom topography, patrol maritime areas, collect water samples, and transport cargo [liu2016usv, bai2022review, Peng2023]. A mission that covers a large area or that measures several points at the same time uses a fleet of USVs, and a ground station (GS) coordinates the vessels of the fleet [kim2026Command, couceiro2026Paradigm].

A digital twin (DT) is a digital representation of a physical object that exchanges data with this object automatically in both directions [kritzinger2018digital]. The DT of a fleet holds one twin per vehicle and updates each twin with the data of its vehicle [yan2026Swarm, zhang2024Formation], so a GS that runs the twins needs the state of every vessel in every period of the twins. The vessels and the GS exchange these states and the decisions of the twins over one wireless network, and the vessels contend for the channel of this network [Shang2025, bianchi2000performance]. A network digital twin (NDT) is a DT whose physical object is a data network [zhou2026ndt, almasan2022network], and an NDT evaluates a configuration of the network before the network applies it [almasan2022network].

The data of one vessel differ in size by orders of magnitude, from the state of the vessel, a few tens of bytes per update, to the frames of a stereo camera, tens of kilobytes per frame [zed]. When the telemetry, the states, and the video of a vessel share one Transmission Control Protocol (TCP) connection, which delivers the bytes in order [rfc9293], and the vessels offer more bytes than the channel carries, the states wait behind the video, and the GS receives the states late or in bursts [icecs2026fleetdt]. The twins of the GS then compute their decisions from late states.

The related works reduce the load of each vehicle at the source [mtowe2025Edge, xu2025Dynamic, kremer2023Teleoperated], reserve bandwidth for each MQTT flow in a wired network [shahri2022Extending], or model wired networks [almasan2022network, rusek2020routenet]. The DBA of passive optical networks and of IEEE~802.16, the controlled access of IEEE~802.11, and the overlay TDMA schemes over commodity IEEE~802.11 divide a shared medium among the stations in each cycle or slot [kramer2002ipact, ieee80216std2017, ieee80211std2020, rao2005overlay, djukic2009soft]. These works leave open the allocation of one shared Wi-Fi channel among the vessels of a fleet in each period, in the order of the data that the GS needs, with a bound on the fraction of the periods in which the GS misses an allocated message for any behavior of the channel. This paper addresses this allocation with three research questions (RQs):

\noindent RQ.I: under which condition on the capacity of the channel does the GS receive, within a period, every message that the allocation admits, and which fraction of the periods can the allocation guarantee for any behavior of the channel?

\noindent RQ.II: how accurately does a prediction from the busy time of the channel and a model of the IEEE~802.11 DCF estimate the capacity of the next period, compared with predictions from the measured capacity alone?

\noindent RQ.III: does the admission of the data of the fleet before each period keep the delivery of the data that the GS needs when the load of the fleet exceeds the capacity of the channel, compared with the prioritization of the IEEE~802.11 MAC, strict priority, and overlay TDMA?

\noindent Section~\ref{subsec:admission} answers RQ.I with Lemma~\ref{lem:delivery} and Theorem~\ref{thm:aci}, and the experiments of Section~\ref{sec:method} test RQ.I and answer RQ.II and RQ.III.

This paper presents Fleet-DT, an NDT that allocates one Wi-Fi channel among the vessels of a fleet in steps of $T_{\mathrm{s}} = 125$~ms by dynamic bandwidth allocation (DBA) above the IEEE~802.11 MAC. In each step, each vessel reports the demand of its data and the bytes that its MQTT bridge still holds; the GS predicts the capacity of the channel from the busy time that the access point measures and from the efficiency of the DCF, corrects the prediction with adaptive conformal inference (ACI) [gibbs2021adaptive], and admits the data of the step in the order of need of its applications within the corrected capacity; and the regulator in the MQTT broker of each vessel releases the admitted data and withholds the other data.

This paper extends the conference paper [icecs2026fleetdt]. With respect to [icecs2026fleetdt], this paper
(i) twins the network of the fleet with the DT prototype (DTP), the DT instance (DTI), the DT aggregate (DTA), and the DT environment (DTE) of Grieves and Vickers [grieves2017], with the state of the network of each step in the DTI and a what-if evaluation of configurations in the DTE (Section~\ref{subsec:dts});
(ii) defines the capacity of a step as a quantity that the access point measures and predicts it with the DCF model of Bianchi and the busy time of the channel (Section~\ref{subsec:pred});
(iii) admits the units of data of every vessel within one airtime budget in the order of need of the GS, with a required set of topics in every step and whole video frames, and proves that the admission returns the lexicographically largest feasible set (Section~\ref{subsec:admission});
(iv) proves that the GS receives every admitted message within the step when the work of the step fits in the realized capacity, and that ACI bounds the long-run fraction of the steps with a missed admitted message for every realization of the channel (Section~\ref{subsec:admission});
(v) states the requirements of a platform and the regulator that releases the admitted data at the pace of the radio (Section~\ref{sec:platform}); and
(vi) evaluates Fleet-DT with the unchanged software on emulated hosts over a real-time ns-3 emulation of the channel, with hypotheses, thresholds, and seeds fixed in a tagged commit before the evaluation (Sections~\ref{sec:method} and~\ref{sec:results}).

Section~\ref{sec:background} presents the concepts of the literature on which Fleet-DT builds and compares Fleet-DT with the related works, Section~\ref{sec:problem} defines the network and states the allocation problem, Section~\ref{sec:dt} defines the model of Fleet-DT, and Section~\ref{sec:platform} states the requirements of a platform. Section~\ref{sec:method} presents the experimental methodology, Section~\ref{sec:results} the results, and Section~\ref{sec:discussion} the discussion. Section~\ref{sec:conclusion} concludes the paper.

\section{Background and Related Work}\label{sec:background}

This section presents the concepts of the literature on which Fleet-DT builds and compares Fleet-DT with the related works in four axes: the DTs of vehicles, fleets, and networks (Section~\ref{subsec:bg_dt}), MQTT and the transport of the data of DTs (Section~\ref{subsec:bg_mqtt}), the congestion of a shared wireless link (Section~\ref{subsec:cong}), and the allocation of a shared medium and the regulation of traffic (Section~\ref{subsec:dba}). Section~\ref{subsec:aci} presents the one-step prediction of a measured series, adaptive conformal inference, and the age of information. Table~\ref{tab:related} compares Fleet-DT with the standards and the research systems closest to it.

% Table I of Fleet-DT v2. Requires \usepackage{booktabs,tabularx}.
% Citation keys follow the [key] convention of the revision files; replace by \cite{key} in the IEEE source.
% Every value comes from the cited work (checked 2026-10-15, see REFS_REPORT.md); n.r. marks a value
% that the cited work does not report or that could not be checked. Vehicle DTs without a channel
% allocation (Kremer, Mtowe, Xu, Yan, Zhang, Baniqued, Hartwell) stay in the text of Section II-A/II-B.
\begin{table*}[t]
\caption{Comparison of Fleet-DT with the standards and the research systems closest to it.}
\label{tab:related}
\centering\footnotesize
\setlength{\tabcolsep}{3pt}
\begin{tabularx}{\textwidth}{@{}p{3.0cm}p{2.3cm}p{2.5cm}XXp{2.6cm}@{}}
\toprule
Work & Layer & Allocation unit & Guarantee & Prediction & Validation \\
\midrule
\multicolumn{6}{@{}l}{\emph{Standards}} \\
IEEE 802.11 EDCA [ieee80211std2020] & MAC & Access category & Priority by AIFS, contention window, and TXOP; no bound & None & Standard \\
IEEE 802.11 HCCA [ieee80211std2020] & MAC & Flow (TSPEC) & Admission by the AP; polled transmissions; delay bound declared in the TSPEC & None & Standard \\
IEEE 802.11ax [ieee80211ax2021, khorov2019tutorial] & MAC and PHY & Resource unit per trigger frame; TWT service period & Uplink schedule of the AP & None & Standard \\
IEEE 802.11be [ieee80211be2024, adame2021tsn] & MAC & Restricted TWT service period & Service periods reserved for latency-sensitive traffic & None & Standard \\
IEEE 802.16 [ieee80216std2017] & MAC & Uplink grant per frame (UL-MAP) & Request--grant per frame & None & Standard \\
ETSI DCC [etsi2018dcc] & Access layer & Transmissions of a station & Channel load limited from the CBR measured over 100~ms & None (reacts to the CBR) & Standard \\
\midrule
\multicolumn{6}{@{}l}{\emph{Research systems: medium allocation and admission control}} \\
Kramer et al.\ [kramer2002ipact], 2002 & MAC (EPON) & Window per ONU per polling cycle & Grant capped at a maximum window (limited service) & None (reported queue) & Simulation \\
Rao and Stoica [rao2005overlay], 2005 & Overlay above the IEEE 802.11 MAC & Equal time slot per node & Slots allocated by a distributed algorithm & None & Simulation and Click testbed \\
Yang and Kravets [yang2005contention], 2005 & Network (ad hoc) & Flow & Admission when the nodes in carrier-sensing range have the consumed bandwidth & Available bandwidth from the idle time (moving average) & Analysis and ns-2 simulation \\
Leith et al.\ [leith2005tcp], 2005 & MAC (EDCA parameters) & Traffic class of the AP and of the stations & TCP upload and download fairness & None & Testbed and NS simulation \\
Sarr et al.\ [sarr2008bandwidth], 2008 & Network (ad hoc) & Flow & Admission when the estimate covers the rate of the flow & Available bandwidth from idle overlap and collision probability & n.r. \\
Djukic and Mohapatra [djukic2009soft], 2009 & Software TDMA over commodity IEEE 802.11 & TDMA slot & Clocks synchronized within microseconds; arbitrary TDMA schedule & None & Multi-hop testbed \\
Winstein et al.\ [winstein2013stochastic], 2013 & Transport (end to end) & Bytes per forecast interval & Bounded risk of a delay beyond a target & Stochastic forecast of the link deliveries & Traces of four LTE and 3G networks \\
Wei et al.\ [wei2013rtwifi], 2013 & MAC (TDMA on the IEEE 802.11 TSF) & TDMA slot & Slots at fixed instants & None & Testbed \\
H{\o}iland-J{\o}rgensen et al.\ [hoiland2017ending], 2017 & MAC queues of the AP (Linux) & Airtime per station & Airtime fairness & None & Testbed \\
Shahri et al.\ [shahri2022Extending], 2022 & Application (MQTT) over SDN & MQTT flow & SDN bandwidth reservation; period and deadline per flow & None & Mininet emulation \\
Jaiswal et al.\ [jaiswal2023wiros], 2023 & Application (ROS 2) and EDCA parameters & ROS 2 topic & Desired QoS in a static traffic scenario & None & n.r. \\
\midrule
\multicolumn{6}{@{}l}{\emph{Research systems: network digital twins and digital twin synchronization}} \\
Rusek et al.\ [rusek2020routenet], 2020 & Model of a wired network & Path & None & Delay, jitter, and loss per path (GNN) & Packet-level simulation (OMNeT++) \\
Almasan et al.\ [almasan2022network], 2022 & NDT of a wired network & Routing configuration & None & Delay per path (RouteNet-E) in what-if analysis & Case study against OMNeT++ \\
Almeida et al.\ [almeida2023position], 2023 & NDT of an IEEE 802.11 network (ns-3) & Link & None & Throughput in ns-3 with a learned propagation loss & ns-3 against testbed \\
Li et al.\ [li2024aoi], 2024 & Edge computing & Twin of an object & $1/4$ approximation (static); competitive ratio (dynamic) & None (AoI computed) & Simulation \\
An et al.\ [an2025enabling], 2025 & Application (MQTT) & Update of a UAV twin & None (AoI measured) & None & Simulation \\
\midrule
Fleet-DT & Application (MQTT brokers) above the IEEE 802.11 EDCA & Admission unit (vessel, topic, output) per step & Failure fraction at most $\beta + \alpha_1/(\gamma T) + (1-\beta)(1 + F_{\mathcal{R}}(T))/T$ for every realization of the channel (Theorem~\ref{thm:aci}) & $\hat{\kappa}(t) = (1 - \hat{b}(t))\,\eta(M+1)$ with ACI on the budget & Emulation (ns-3 real time, KVM hosts) \\
\bottomrule
\end{tabularx}
\end{table*}


\subsection{Digital Twins of Vehicles, Fleets, and Networks}\label{subsec:bg_dt}
Kritzinger et al.\ classify virtual representations by the automation of the data flows between a physical object and its representation [kritzinger2018digital]. A digital model exchanges data with the physical object only through manual flows, a digital shadow (DS) receives the data of the physical object automatically, and a DT exchanges data with the physical object automatically in both directions. Grieves and Vickers define the DTP, the DTI, the DTA, and the digital twin environment (DTE) [grieves2017]. The DTP describes a prototypical physical artifact and contains the information necessary to produce a physical version of the artifact. The DTI describes one physical product and remains linked to that product throughout the life of the product. The DTA is the aggregation of the DTIs. The DTE is the integrated multi-domain physics application space in which applications operate on the DTs [grieves2017]. Section~\ref{subsec:dts} defines how Fleet-DT adopts these constructs.

Peng et al.\ [Peng2023] and Wang et al.\ [wang2026Marine] surveyed DTs for unmanned surface and underwater vehicles, and Wang et al.\ classified one of eleven marine robotic cases as a bidirectional DT. Most single-vehicle works implement a DS [kritzinger2018digital]. Bhagavathi et al.\ [Bhagavathi2023] fed sensor data of the Otter autonomous surface vehicle into fault diagnosis models, Berg et al.\ [berg2025Syncing] synchronized the twin of an autonomous surface vehicle by parameter identification with deep reinforcement learning, and Xu et al.\ [xu2025Dynamic] sampled the state of a ship on events to reduce the ship-to-shore traffic. These works couple one twin to one vehicle, leave the transport unspecified, and state no deadline for the update of the twin.

Gervasi et al.\ [gervasi2026Systematic] classified systems of DTs into DTAs, DT networks, federated DTs, composite DTs, and DT ecosystems. Yan et al.\ [yan2026Swarm] and Zhang et al.\ [zhang2024Formation] assigned one twin to each autonomous underwater vehicle of a swarm or to each remotely operated vehicle of a formation and updated the twins periodically. Kim et al.\ [kim2026Command] validated the command-and-control system of five USVs in a DT before the field operation, Wang et al.\ [wang2026Optimizing] trained a cooperative cruising policy for surface vehicles in a DT of the fleet, and Couceiro et al.\ [couceiro2026Paradigm] ran a Unity twin of an aquatic robot swarm with ROS~2 over a Wi-Fi mobile ad hoc network. Outside the marine domain, Baniqued et al.\ [Baniqued2024], Mtowe et al.\ [mtowe2025Edge], and Hartwell et al.\ [Hartwell2024] twinned robot and aero-engine fleets. These twins serve as validation or training environments or as control models. Fleet-DT twins the shared Wi-Fi network of a fleet, with the state of the network of each step in the DTI and the DTIs of the past operations in the DTA.

Zhou et al.\ describe the network digital twin (NDT) as a digital representation whose physical counterpart is a data network, with four elements [zhou2026ndt]. The data hold the historical and real-time state of the network, the models emulate changes in the configuration, the state, or the use of the network elements, the mapping links the NDT to the network in both directions in real time, and the interfaces connect the NDT to the network and to the applications [zhou2026ndt]. Almasan et al.\ represent the network as a graph, use graph neural networks (GNNs) as the models of the NDT, and apply the NDT to what-if analysis, optimization, and troubleshooting [almasan2022network]. RouteNet, the GNN of Rusek et al., takes the topology, the routing, and the traffic matrix of a network and predicts the mean delay, the jitter, and the loss of each path, and the authors trained and tested RouteNet on packet-level simulations of wired topologies with 14 to 50 nodes [rusek2020routenet]. These NDTs model wired networks with routing over several hops. Almeida et al.\ built a DT of a Wi-Fi network in the ns-3 network simulator with a propagation loss model that the authors trained on testbed traces, and this DT runs without a link to the network in operation [almeida2023position]. Fleet-DT twins one Wi-Fi network in operation, in which each vessel reaches the GS in one hop, and predicts, in each step, the capacity of the shared channel from the measurements of this network.

Wu et al.\ survey digital twin networks, in which the twins of physical objects exchange data with these objects over a network [wu2021digital]. Kapteyn et al.\ formulate a DT as a probabilistic graphical model that couples the physical state, the observed data, the digital state, and the control actions over time, and update the digital state of an unmanned aerial vehicle in flight by Bayesian inference [kapteyn2021probabilistic]. The National Academies require continual verification, validation, and uncertainty quantification of a DT for the purpose of the DT [nasem2024foundational]. Fleet-DT states the fidelity of the DT of the network with the metrics of Definition~\ref{def:fidelity} and measures each metric against the ground truth of the emulator (Section~\ref{sec:method}).

\subsection{MQTT and the Transport of the Data of Digital Twins}\label{subsec:bg_mqtt}
In the publish/subscribe protocol MQTT, a client publishes messages on topics to a broker, and the broker delivers each message to the clients that subscribe to the topic of the message [oasis2014mqtt]. MQTT requires a network connection that delivers the bytes in order and without loss, and MQTT implementations use TCP for this connection [oasis2014mqtt]. An MQTT bridge connects two brokers over one client connection and forwards the messages of selected topics between the two brokers, and the configuration of the bridge gives the direction and the QoS level of each forwarded topic [eclipse2026mosquittoconf].

MQTT defines three QoS levels for the delivery of a message from a sender to a receiver [oasis2014mqtt]. With QoS~0, the sender sends the PUBLISH packet once, and the receiver sends no acknowledgment (at most once). With QoS~1, the receiver acknowledges the PUBLISH packet with a PUBACK packet, and the sender keeps the message until the PUBACK packet arrives (at least once). With QoS~2, the sender and the receiver exchange the PUBLISH, PUBREC, PUBREL, and PUBCOMP packets (exactly once). When a client reconnects to a broker with a persistent session (CleanSession flag set to 0), the client and the broker each resend the PUBLISH packets of QoS~1 and~2 that the other side did not acknowledge [oasis2014mqtt]. Each QoS level governs one hop, from a client to a broker or from a broker to a client [oasis2014mqtt].

Zhang et al.\ [Zhang2024] compared DDS, MQTT, and Zenoh under ROS~2 over Ethernet, Wi-Fi, and 4G and measured the lowest delay with Zenoh on the wireless links. Freitas et al.\ [freitas2026Distribution] evaluated MQTT, AMQP, Kafka, DDS, and Zenoh in a DT synchronization architecture and reported that DDS collapsed with 108 nodes under stress, while MQTT and AMQP synchronized states of up to 2~MiB on a degraded network. An et al.\ [an2025enabling] proposed a centralized MQTT framework that updates the DTs of a simulated UAV network with lower latency, jitter, and age of information than MQTT, TCP sockets, and CoAP. Lima et al.\ [lima2025Bridge] showed that the number of MQTT bridges, the packet size, and the topic name affect latency and reliability, and Li et al.\ [li2023Remote] carried sensor data of the research vessel Gunnerus over MQTT to a remote control center.

For video, Herrero [herrero:2023] transported media packets inside MQTT messages, and Kremer et al.\ [kremer2023Teleoperated] and Mtowe et al.\ [mtowe2025Edge] replaced the video stream with a twin of the environment or with edge features to save bandwidth; RTP [rfc3550] carries media in its own session with RTCP feedback. For QoS, Shahri et al.\ [shahri2022Extending] let each MQTT flow declare a period and a deadline and reserved bandwidth through SDN, Palmese et al.\ [palmese2022Adaptive] adjusted the QoS level of each MQTT-SN node at the broker from delay and packet error, and the DDS DEADLINE policy [omg2015Dds] signals missed periodic updates. Jaiswal et al.\ [jaiswal2023wiros] scheduled the flows of the ROS~2 topics of one robot and configured the IEEE~802.11 contention parameters of its Wi-Fi interface for the latency-critical topics. Fleet-DT carries the telemetry, the states, and the video of each vessel over the data bridge of the vessel and admits, in each step, the data of every vessel within the capacity that the DT of the network predicts for the step.

\subsection{Congestion of a Shared Wireless Link and Time-Triggered Communication}\label{subsec:cong}
In IEEE~802.11, the access point and its stations contend for one channel with the distributed coordination function (DCF), so the traffic of both directions shares the capacity of the link [bianchi2000performance]. Bianchi defines the saturation throughput of the DCF as the throughput of the channel when the transmission queue of every station always holds a packet, and shows that the saturation throughput depends on the number of stations that contend for the channel, on the size of the frames, and on the parameters of the physical layer [bianchi2000performance]. Heusse et al.\ [heusse2003performance] show that a station that transmits at a lower bit rate reduces the throughput of every station of the cell and name this effect the performance anomaly of IEEE~802.11. TCP delivers one ordered byte stream per connection [rfc9293], so a segment that waits in the queue delays every later byte of the same connection.

Bianchi models the backoff of each of $n$ saturated stations with a Markov chain and obtains the probability $\tau$ that a station transmits in a slot and the probability $p_{\mathrm{c}}$ that a transmitted frame collides from
\begin{equation}
\tau = \frac{2\,(1 - 2p_{\mathrm{c}})}{(1 - 2p_{\mathrm{c}})(\mathrm{CW}_{\min} + 2) + p_{\mathrm{c}} (\mathrm{CW}_{\min} + 1) \big(1 - (2p_{\mathrm{c}})^{m_{\mathrm{b}}}\big)}, \qquad p_{\mathrm{c}} = 1 - (1 - \tau)^{n-1},
\label{eq:bianchi}
\end{equation}
where $m_{\mathrm{b}}$ is the maximum backoff stage [bianchi2000performance]. With $P_{\mathrm{tr}} = 1 - (1 - \tau)^n$, the probability that a slot holds a transmission, and $P_{\mathrm{s}} = n\tau(1-\tau)^{n-1}/P_{\mathrm{tr}}$, the probability that this transmission succeeds, the fraction of time in which the channel carries successful transmissions is
\begin{equation}
\eta(n) = \frac{P_{\mathrm{tr}} P_{\mathrm{s}}\, T_{\mathrm{succ}}}{(1 - P_{\mathrm{tr}})\,T_{\mathrm{slot}} + P_{\mathrm{tr}} P_{\mathrm{s}}\, T_{\mathrm{succ}} + P_{\mathrm{tr}}(1 - P_{\mathrm{s}})\, T_{\mathrm{coll}}},
\label{eq:eff}
\end{equation}
where $T_{\mathrm{slot}}$ is the slot time, $T_{\mathrm{succ}}$ is the duration of a successful transmission (PHY header, MAC header and payload, SIFS, acknowledgment, and DIFS), and $T_{\mathrm{coll}}$ is the duration of a collision (PHY header, MAC header and payload, and DIFS) [bianchi2000performance]. The saturation throughput of Bianchi is~\eqref{eq:eff} with the transmission time of the mean payload in place of $T_{\mathrm{succ}}$ in the numerator. With EDCA, AIFS of the access category takes the place of DIFS and the contention window of the access category takes the place of $\mathrm{CW}_{\min}$ [ieee80211std2020]. The ns-3 Model Library compares the throughput of the ns-3 Wi-Fi model with the model of Bianchi for IEEE~802.11a, b, and g [nsnam2020wifi].

IEEE~802.11 adds to the DCF the hybrid coordination function (HCF) [ieee80211std2020]. With EDCA, each station keeps one queue per access category, and each access category contends for the channel with its own interframe space and contention window. With HCCA, which IEEE~802.11 defines as optional, a station declares the mean data rate, the maximum burst size, and the delay bound of a flow in a TSPEC, the access point admits or rejects the flow (admission control), and the access point polls the stations of the admitted flows [ieee80211std2020]. In vehicular networks, the ETSI DCC adapts the transmissions of each station to the channel busy ratio (CBR), the fraction of time in which the channel is busy with transmissions, which each station measures over intervals of 100~ms [etsi2018dcc]. The available bandwidth of a link in an interval is the capacity of the link that the current traffic leaves unused in the interval [prasad2003bandwidth].



IEEE~802.11 maps the user priority (UP) of a frame to the access category, with UP~6 and~7 in AC\_VO and UP~0 and~3 in AC\_BE, and RFC~8325 maps the differentiated services code point (DSCP) of an IP packet to the UP [ieee80211std2020, rfc8325]. Leith et al.\ show that, under the DCF, the AP of an infrastructure WLAN receives about the same share of the transmission opportunities as one station, although the AP sends the TCP acknowledgments of every upload and the data of every download, and that EDCA parameters that give the acknowledgments at the AP a shorter AIFS and a smaller $\mathrm{CW}_{\min}$ and the download data at the AP a TXOP of several packets restore the fairness between TCP uploads and downloads [leith2005tcp].

In the time-triggered architecture, the nodes of a distributed system act at instants of a global time base that a clock synchronization establishes, so every node knows in advance the instants at which each message is sent [kopetz2003time]. A firm deadline makes a result useless after the deadline [kopetz2011real].

\subsection{Allocation of a Shared Medium and Traffic Regulation}\label{subsec:dba}
In DBA, a controller divides the capacity of a shared medium among the stations in each cycle from the demand that each station reports. In EPONs, the optical line terminal grants each optical network unit a transmission window in each polling cycle from the queue length that the unit reports, and the limited service of IPACT caps each grant at a maximum window [kramer2002ipact]. In IEEE~802.16, the subscriber stations send bandwidth requests, and the base station grants the uplink bandwidth of each frame in the uplink map [ieee80216std2017]. Both schemes repeat a request--grant cycle, in which the stations report their demand, the controller computes the grants, and the stations transmit within their grants.

Rao and Stoica divide the time into equal slots with loosely synchronized clocks above the IEEE~802.11 MAC and allocate the slots among the nodes with a distributed algorithm [rao2005overlay]. Djukic and Mohapatra synchronize commodity IEEE~802.11 nodes within microseconds and run TDMA schedules over several hops [djukic2009soft]. Wei et al.\ build RT-WiFi, a TDMA MAC on the timing synchronization function of IEEE~802.11, for wireless control loops [wei2013rtwifi]. IEEE~802.11ax lets the AP schedule uplink transmissions on resource units with trigger frames and the service periods of the stations with target wake time (TWT) [ieee80211ax2021, khorov2019tutorial], and IEEE~802.11be adds restricted TWT for latency-sensitive traffic [ieee80211be2024, adame2021tsn]. These mechanisms give slots or resource units to nodes, and each node selects the data of its slot. Fleet-DT admits the data of every vessel within one airtime budget per step in the order of need of the GS, on the EDCA of radios without trigger-based access (Section~\ref{sec:platform}), and Section~\ref{sec:method} compares Fleet-DT with an overlay TDMA and with the slotted variant of Definition~\ref{def:tdma}.

Network calculus bounds the traffic of a flow with an arrival curve, the function that gives, for each duration, the largest number of bytes that the flow sends in an interval of this duration, and a token bucket with a rate and a bucket size gives the arrival curve that starts at the bucket size and grows with the rate [cruz1991calculus, leboudec2001network]. A regulator delays the data of a flow until its output respects an arrival curve, a policer discards or marks the data that exceed an arrival curve, and a strict-priority scheduler serves a class only when the classes of higher priority hold no data [leboudec2001network]. Fleet-DT applies the request--grant cycle above the IEEE~802.11 MAC, to the units of data of the vessels in the order of need of the GS, and regulates the release of the admitted units in each vessel (Section~\ref{sec:dt}).

\subsection{One-Step Prediction and Adaptive Conformal Inference}\label{subsec:aci}
Exponential smoothing with the gain $\lambda \in (0, 1]$ predicts the next value of a measured series $y(t)$ with $\hat{y}(t+1) = (1 - \lambda)\,\hat{y}(t) + \lambda\, y(t)$, and, for a random walk observed with additive noise, a gain that depends on the ratio of the variance of the increments to the variance of the noise minimizes the mean squared error of the one-step prediction [harvey1989forecasting]. TCP predicts the round-trip time of a connection with exponential smoothing and the gain $1/8$ [rfc6298].

Conformal prediction turns a point predictor into a prediction set: with a nonconformity score $S_t$ and the empirical quantile $\hat{Q}_t$ of the past scores, the set $\hat{C}_t(\alpha) = \{y : S_t(y) \leq \hat{Q}_t(1 - \alpha)\}$ covers the next observation with probability at least $1 - \alpha$ when the observations are exchangeable [gibbs2021adaptive]. Adaptive conformal inference (ACI) requires no exchangeability of the observations. With $\mathrm{err}_t = 1$ when the observation of step $t$ falls outside $\hat{C}_t(\alpha_t)$ and $\mathrm{err}_t = 0$ otherwise, ACI updates the level with
\begin{equation}
\alpha_{t+1} = \alpha_t + \gamma\,(\beta - \mathrm{err}_t),
\label{eq:aci_bg}
\end{equation}
where $\beta$ is the target miscoverage and $\gamma > 0$ is the step. The conventions $\hat{Q}_t(x) = \infty$ for $x > 1$ and $\hat{Q}_t(x) = -\infty$ for $x < 0$ keep $\alpha_t \in [-\gamma, 1 + \gamma]$, and summing~\eqref{eq:aci_bg} gives $\big|\frac{1}{T}\sum_{t=1}^{T} \mathrm{err}_t - \beta\big| \leq (\max\{\alpha_1, 1 - \alpha_1\} + \gamma)/(\gamma T)$ for every sequence of observations [gibbs2021adaptive].

The age of information (AoI) of a status at a monitor is the time elapsed since the release of the newest status that the monitor has received [kaul2012real].
Kaul et al.\ show for a first-come first-served queue that the average AoI grows when the update rate approaches zero or the service rate, so an intermediate update rate minimizes the average AoI [kaul2012real]. Yates et al.\ survey the AoI of queues, of scheduled sources, and of sources that share a channel [yates2021age]. Li et al.\ place the twins of mobile objects on edge servers, serve each query from the twin with the lowest AoI, and give a $1/4$ approximation for the static placement [li2024aoi]. These works minimize an average AoI over the updates of the sources. Fleet-DT admits the state of every vessel in every step and bounds the AoI of the state at the GS in every step that meets the hypothesis of Lemma~\ref{lem:delivery} (Corollary~\ref{cor:state}).

\section{Problem Statement}\label{sec:problem}

\subsection{Network}\label{subsec:network}

\begin{definition}[Fleet]\label{def:fleet}
The fleet comprises the vessels $i \in \{1, \dots, M\}$, with $M \geq 2$, and one ground station (GS) with the index $0$. The vessels and the GS are hosts of one IEEE~802.11 basic service set: the radio of the GS is the access point (AP), and the radio of each vessel is a station associated with the AP. The uplink carries the messages from the vessels to the GS, and the downlink carries the messages in the opposite direction.
\end{definition}

Each vessel and the GS run one MQTT broker. The broker of the GS opens one data bridge to the broker of each vessel over one TCP connection, and the data bridge carries the $L$ uplink topics of the vessel. Every vessel has the same $L$ uplink topics.

\begin{definition}[Uplink topic]\label{def:topic}
Uplink topic $\ell \in \{1, \dots, L\}$ has one source in each vessel. The source releases $r_\ell$ outputs per second with a release jitter of at most $J_\ell$, and each output comprises at most $n_\ell$ messages of at most $s_\ell$ bytes of content. Topic $\ell$ is of the kind \emph{stream} when the consumer of topic $\ell$ at the GS uses each message alone, and of the kind \emph{output} when the consumer uses only complete outputs, each output usable without the other outputs.
\end{definition}

Table~\ref{tab:topics} gives $r_\ell$, $J_\ell$, $n_\ell$, $s_\ell$, and the kind of the $L = 7$ uplink topics of the experiments. The six telemetry topics are streams with $n_\ell = 1$, and each output of the video topic is one intra-coded image, which its source splits into at most $n_\ell$ messages.

\begin{definition}[Required set]\label{def:required}
The required set $\mathcal{R} \subseteq \{1, \dots, L\}$ holds the uplink topics whose consumers at the GS require the messages of every vessel in every step. The messages of the topics of $\mathcal{R}$ cross the data bridge with MQTT QoS~2, and the messages of the other topics cross the data bridge with MQTT QoS~0.
\end{definition}

The applications of the vessels and of the GS fix $\mathcal{R}$ and the QoS level of $\mathcal{R}$, and Fleet-DT receives both as data. The downlink topic \texttt{goals} carries the goals of each vessel in the operation, with a rate and a size bound that the DTP declares (Section~\ref{subsec:dts}).

\subsection{Steps and Airtime}\label{subsec:steps}

\begin{definition}[Step]\label{def:step}
Step $t \in \{1, 2, \dots\}$ is the interval $[t_0 + (t-1)\,T_{\mathrm{s}},\ t_0 + t\,T_{\mathrm{s}})$ of the clock of the GS, where $t_0$ is the start of the operation and $T_{\mathrm{s}}$ is the duration of a step.
\end{definition}

The source of the topic \texttt{state} releases $r_{\mathrm{state}} = 8$ outputs per second in the experiments, and $T_{\mathrm{s}} = 1/r_{\mathrm{state}} = 125$~ms, so each step carries one state of each vessel. Section~\ref{subsec:window} gives the overhead of the phases and the delay of a state as functions of $T_{\mathrm{s}}$, and Section~\ref{sec:method} varies $T_{\mathrm{s}}$.

Each IP packet crosses the channel in one IEEE~802.11 frame, and the receiver of the frame answers with a MAC acknowledgment. The airtime $T_{\mathrm{succ}}(s, R)$ of a successful exchange of an IP packet of $s$ bytes at the PHY rate $R$ comprises the PHY preamble and header, the MAC header and the $s$ bytes at the rate $R$, SIFS, the MAC acknowledgment at the control response rate of $R$, and the arbitration interframe space (AIFS) of the access category of the frame [bianchi2000performance, ieee80211std2020]. The backoff slots and the collisions stay outside $T_{\mathrm{succ}}$.

\begin{definition}[Airtime cost]\label{def:cost}
The airtime cost $a_i(x, n, q)$ of $n$ MQTT messages of QoS $q$ with $x$ bytes of content in total, sent by vessel $i$ over its data bridge, is an upper bound on the sum of $T_{\mathrm{succ}}$ over the frames of the TCP segments that carry these messages, the frames of the TCP acknowledgments of these segments, and the frames of the MQTT control packets of QoS $q$ with their TCP acknowledgments, at the PHY rate $R_i$ for the frames of vessel $i$ and at the PHY rate $R_{0,i}$ for the frames of the AP to vessel $i$.
\end{definition}

The bound of Definition~\ref{def:cost} holds when the PHY rates do not fall below $R_i$ and $R_{0,i}$ and TCP retransmits no segment; Section~\ref{subsec:costplat} gives its closed form.

\subsection{Allocation Problem}\label{subsec:alloc}

The vessels and the GS contend for the channel with the DCF (Section~\ref{subsec:cong}), so the frames of one vessel in a step reduce the airtime that the channel gives to the other vessels in the step. The TCP connection of the data bridge of each vessel sends the bytes of the $L$ uplink topics in the order of their arrival at the socket (Section~\ref{subsec:cong}). When the vessels offer in a step more frames than the channel carries in the step, the excess bytes wait in the send buffers and in the radio queues, and the messages that the GS needs wait behind the other messages (head-of-line blocking).

\begin{definition}[Admission unit]\label{def:unit}
An admission unit of step $t$ is a triple $u = (i, \ell, k)$. For a stream topic $\ell$, the unit $(i, \ell, 1)$ holds the messages of topic $\ell$ of vessel $i$ in step $t$. For an output topic $\ell$, the unit $(i, \ell, k)$ holds the messages of output $k$ of topic $\ell$ of vessel $i$ in step $t$. The demand $(x_u(t), n_u(t))$ of unit $u$ bounds the bytes of content and the number of the messages of $u$, and the cost of $u$ is $c_u(t) = a_i(x_u(t), n_u(t), q_\ell)$, where $q_\ell$ is the QoS of topic $\ell$.
\end{definition}

In each step $t$, the applications of the GS provide, for each vessel $i$, the set $\mathcal{N}_i(t) \supseteq \mathcal{R}$ of the uplink topics that these applications need from vessel $i$ in step $t$, and the order $\sigma_t$, a permutation of $\{1, \dots, L\}$ that ranks the uplink topics by need, with the topics of $\mathcal{R}$ first. Section~\ref{subsec:admission} extends $\sigma_t$ to an order $\pi_t$ of the units of step $t$. A set $A$ of units is lexicographically larger than a set $A'$ in the order $\pi_t$ when the first unit of $\pi_t$ that belongs to exactly one of the two sets belongs to $A$.

The data window $W(t)$ is the part of step $t$ in which the vessels send the messages of step $t$, with the start $w(t)$ and the duration $T_W$ (Section~\ref{subsec:window}). The fleet sends in $W(t)$ the frames of its data bridges and of the bytes that the data bridges hold from earlier steps. A data bridge holds bytes at the end of $W(t)$ when the AP records a frame of the data bridge from the end of $W(t)$ to the end of step $t$ or the vessel reports held bytes in step $t+1$ (Section~\ref{subsec:protocol}). Let $e(t)$ be the end of $W(t)$ when a data bridge holds bytes at the end of $W(t)$, the end of the last successful exchange of a frame of a data bridge in $W(t)$ otherwise, and $w(t)$ when the data bridges exchange no frame in $W(t)$. Let $T_{\mathrm{fleet}}(I)$ be the sum of $T_{\mathrm{succ}}$ over the successful exchanges of the frames of the data bridges in the interval $I$, where an exchange that crosses a boundary of $I$ counts in proportion to its overlap with $I$.

\begin{definition}[Capacity]\label{def:capacity}
The capacity of step $t$ with $e(t) > w(t)$ is $\kappa(t) = T_{\mathrm{fleet}}([w(t), e(t)))/(e(t) - w(t))$, the fraction of the busy period of the fleet in $W(t)$ in which the channel carries successful exchanges of the data bridges.
\end{definition}

The AP exchanges every frame of the data bridges, so the GS obtains $e(t)$ and $\kappa(t)$ in phase~1 of step $t+1$ for every step in which the fleet sends a frame in $W(t)$. The capacity $\kappa(t)$ accounts for the backoff slots, the collisions, the frames of the transmitters outside the fleet and of the traffic outside the data bridges, and the idle time of a TCP connection that waits for an acknowledgment or a retransmission timeout, and $\kappa(t)$ depends on the units that the GS admits.

\begin{problem}[Allocation of the channel]\label{prob:alloc}
In each step $t$, before $W(t)$, choose the airtime budget $K(t)$ and the set $A(t)$ of admitted units such that $A(t)$ holds the units of the topics of $\mathcal{R}$ and is the lexicographically largest set in the order $\pi_t$ among the sets that hold these units and satisfy
\begin{equation}
\sum_{u \in A(t)} c_u(t) \leq K(t),
\label{eq:feas}
\end{equation}
or holds the units of $\mathcal{R}$ alone when no such set exists, and such that, for a target $\beta \in (0, 1)$ and for every realization of the channel, the fraction of the steps $t \leq T$ in which the GS misses within step $t$ a message of a unit of $A(t)$ exceeds $\beta$ by at most a term that vanishes with $T$ and a term that counts the steps in which the channel does not carry the units of $\mathcal{R}$ alone.
\end{problem}

The GS decides $K(t)$ and $A(t)$ before $W(t)$, so the GS knows $\kappa(t)$ only after step $t$. Section~\ref{sec:dt} chooses $K(t)$ from a prediction of $\kappa(t)$ corrected by adaptive conformal inference and solves Problem~\ref{prob:alloc}.

% Tables of Fleet-DT v2 referenced by Sections III, VI.

\begin{table}[t]
\caption{Uplink topics of the experiments (DTP). Each stream output holds one message.}\label{tab:topics}
\centering\footnotesize
\begin{tabular}{@{}lrrrrll@{}}
\toprule
Topic & $r_\ell$ (1/s) & $J_\ell$ (ms) & $n_\ell$ & $s_\ell$ (B) & Kind & QoS \\
\midrule
\texttt{state} ($\mathcal{R}$) & 8 & 2 & 1 & 48 & stream & 2 \\
\texttt{imu} & 32 & 2 & 1 & 48 & stream & 0 \\
\texttt{gps} & 8 & 2 & 1 & 80 & stream & 0 \\
\texttt{baro} & 8 & 2 & 1 & 32 & stream & 0 \\
\texttt{power} & 8 & 2 & 1 & 40 & stream & 0 \\
\texttt{mag} & 16 & 2 & 1 & 40 & stream & 0 \\
\texttt{video} & 15$^{\ast}$ & 10 & 16--23$^{\ast}$ & 1400 & output & 0 \\
\bottomrule
\end{tabular}

\smallskip
\raggedright\footnotesize $^{\ast}$The rate and $n_\ell$ of the video topic follow the offered-load ratio $\rho$ of each run (Section~\ref{subsec:scen}).
\end{table}

\begin{table}[t]
\caption{Parameters of the experiments.}\label{tab:param}
\centering\footnotesize
\begin{tabular}{@{}lll@{}}
\toprule
Group & Parameter & Value \\
\midrule
Step & $T_{\mathrm{s}}$, $T_1$, $T_2$ & 125, 20, 10~ms \\
 & $\varepsilon$, $d_{\mathrm{rx}}$, $g_1$ & 2, 2, 4~ms \\
 & $T_W$ & 89~ms \\
Regulator & $T_{\mathrm{tick}}$, $d_{\mathrm{m}}$ & 1, 1~ms \\
 & $L_{\mathrm{w}}$, $B_{\mathrm{cap}}$ & 16~KiB, 1~MiB \\
AS & $\beta$, $\gamma$, $\alpha_1$ & 0.05, 0.05, 0.5 \\
 & $\lambda$, $T_{\min}$ & 0.75, 5~ms \\
PHY and MAC & Standard, rate & IEEE~802.11g ERP-OFDM, 24~Mb/s$^{\dagger}$ \\
 & Basic rates & 6, 12, 24~Mb/s \\
 & Slot, SIFS & 9, 10~$\mu$s \\
 & AIFSN, $\mathrm{CW}_{\min}$ (stations) & 3, 15 \\
 & AIFSN, $\mathrm{CW}_{\min}$ (AP, own AC\_BE) & 1, 7 \\
 & $\mathrm{CW}_{\max}$, $m_{\mathrm{b}}$ & 1023, 6 \\
Headers & $h_{\mathrm{MAC}}$, $h_{\mathrm{T}}$, $h_{\mathrm{A}}$ & 38, 80, 4~B \\
TCP & MSS, \texttt{rto\_min} & 1448~B, 20~ms \\
\bottomrule
\end{tabular}

\smallskip
\raggedright\footnotesize $^{\dagger}$S5 uses IEEE~802.11a OFDM rates with Minstrel (Section~\ref{subsec:scen}).
\end{table}

\begin{table}[t]
\caption{Validation of the testbed.}\label{tab:valid}
\centering\footnotesize
\begin{tabular}{@{}p{3.6cm}p{4.2cm}@{}}
\toprule
Check & Result \\
\midrule
Saturation throughput, $n = 1, \dots, 16$, 24~Mb/s & within 0.93\% of the ns-3 reference and 1.52\% of~\eqref{eq:eff} (DCF); within 1.6\% of~\eqref{eq:eff} with the EDCA parameters of Table~\ref{tab:param} \\
$\kappa(t)$ and $b(t)$ of the AS against the ground truth (E0) & 99th percentile of the absolute difference $1.2 \times 10^{-4}$ ($\kappa$) and $5.5 \times 10^{-3}$ ($b$) \\
Lag of the real-time scheduler (E0, $M = 4$) & 99th percentile 0.17--0.25~ms (S1, S3--S5) and 0.34--0.54~ms (S2) \\
Clock offset of the vessels after convergence & 30 to 150~$\mu$s, no sample above $\varepsilon$ \\
Monitor path to the broker of the GS & 104 to 119~$\mu$s at the median \\
Hidden vessels (S4) & acknowledged fraction of the frames of the hidden pair 0.52 against 0.9999 without hiding \\
\bottomrule
\end{tabular}
\end{table}

\begin{figure}[t]\centering\resizebox{\columnwidth}{!}{\input{fig_step}}\caption{Step $t$ of Fleet-DT on the clock of the GS: the phases of the GS and of vessel $i$, the demands of the vessels at their offsets in phase~1, the tables of phase~2, the quiet interval $Q(t)$, the data window $W(t)$, and the end $e(t)$ of the busy period of the fleet.}\label{fig:step}\end{figure}
\section{The Digital Twin Model}\label{sec:dt}

The DBA of Fleet-DT solves Problem~\ref{prob:alloc} in each step $t$ with four activities: (1) each vessel sends to the GS the demand of its units, the bytes that its data bridge holds, and the report of step $t-1$; (2) the AS predicts the capacity $\kappa(t)$ and computes the budget $K(t)$ from the prediction and from the level of adaptive conformal inference (ACI); (3) the AS admits the units of step $t$ in the order $\pi_t$ within $K(t)$ and sends to each vessel its table; and (4) the vessel regulator of each vessel releases the admitted units to the data bridge and withholds the other units. Sections~\ref{subsec:window} to~\ref{subsec:admission} define these activities, and Section~\ref{subsec:dts} maps them to the constructs of the DT.

\subsection{Phases and Data Window}\label{subsec:window}

\begin{definition}[Phases of a step]\label{def:phases}
Step $t$ comprises phase~1 (request), phase~2 (grant), and phase~3 (data), in this order, with the durations $T_1 > 0$, $T_2 > 0$, and $T_3 = T_{\mathrm{s}} - T_1 - T_2 > 0$. Phase~1 holds the activity (1), phase~2 holds the activities (2) and (3), and phase~3 holds the activity (4). Each host starts each phase on its own clock.
\end{definition}

\begin{assumption}[Clock offset]\label{as:clock}
In every step $t$, the offset $\theta_i(t)$ of the clock of each vessel $i$ to the clock of the GS satisfies $|\theta_i(t)| \leq \varepsilon$.
\end{assumption}

\begin{definition}[Data window]\label{def:window}
The data window of step $t$ is the interval $W(t) = [w(t),\ t_0 + t\,T_{\mathrm{s}} - g)$ of the clock of the GS, with $w(t) = t_0 + (t-1)\,T_{\mathrm{s}} + T_1 + T_2 + \varepsilon$, $g = \varepsilon + d_{\mathrm{rx}}$, and the duration $T_W = T_3 - 2\varepsilon - d_{\mathrm{rx}}$, where $d_{\mathrm{rx}}$ bounds the time from the end of an exchange at the radio of the GS to the delivery of its message by the broker of the GS.
\end{definition}

Fig.~\ref{fig:step} shows the phases, the data window, and the quiet interval $Q(t)$ (Section~\ref{subsec:pred}) of step $t$. By Assumption~\ref{as:clock}, each vessel starts phase~3 of step $t$ no later than $w(t)$ and ends step $t$ no earlier than the end of $W(t)$ on the clock of the GS. A message of step $t$ reaches its regulator no earlier than the start of phase~3 of step $t-1$ on the clock of its vessel (Section~\ref{subsec:protocol}), so the GS accepts a message of step $t$ within step $t$ at most $T_{\mathrm{s}} + T_3 + \varepsilon$ after its release. The fraction $(T_1 + T_2 + 2\varepsilon + d_{\mathrm{rx}})/T_{\mathrm{s}}$ of each step carries no data of the step, so a shorter step lowers this delay and raises this overhead.

The broker of the GS also opens one control bridge to the broker of each vessel over a second TCP connection. The control bridge carries the topic \texttt{demand} (uplink) and the topics \texttt{levels} and \texttt{goals} (downlink). The hosts give the packets of the control bridge the IEEE~802.11 user priority~6 (AC\_VO) and the packets of the data bridge the user priority~0 (AC\_BE), so a table of phase~1 or~2 waits behind no frame of a data bridge in the queues of the radios and contends for the channel with the EDCA parameters of AC\_VO (Section~\ref{subsec:cong}).

The set $\mathcal{M}(t)$ holds the $m(t)$ vessels whose demand of step $t$ reaches the GS in phase~1 of step $t$; the GS accepts a demand of step $t$ that arrives up to $\varepsilon$ before the start of step $t$.

\subsection{Request--Grant Protocol and Vessel Regulator}\label{subsec:protocol}

Let $v_i(t)$ be the start of phase~3 of step $t$ on the clock of vessel $i$. The collection window of step $t$ of vessel $i$ is $C_i(t) = [v_i(t-1), v_i(t))$, and the messages of step $t$ of vessel $i$ are the messages of the stream topics that reach the regulator of vessel $i$ in $C_i(t)$ and the messages of the outputs whose last message reaches the regulator in $C_i(t)$. At the start of phase~1 of step $t$, the regulator holds the messages of step $t$ that arrived in the first $T_3$ of $C_i(t)$.

\begin{definition}[Demand, held bytes, and report]\label{def:demand}
In phase~1 of step $t$, at the offset $i\,(T_1 - g_1)/(M+1)$ from the start of phase~1, with $g_1 = \varepsilon + 2T_{\mathrm{tick}}$, which spreads the demands of the $M$ vessels over phase~1 and leaves $g_1$ for the last demand to reach the AS, the DTI program of vessel $i$ publishes on the topic \texttt{demand}:
(a) for each unit $u$ of vessel $i$, the demand $x_u(t) = x^{\mathrm{h}}_u + x^{\mathrm{a}}_u$, $n_u(t) = n^{\mathrm{h}}_u + n^{\mathrm{a}}_u$, where $(x^{\mathrm{h}}_u, n^{\mathrm{h}}_u)$ counts the bytes and the messages of $u$ that the regulator holds and $(x^{\mathrm{a}}_u, n^{\mathrm{a}}_u)$ bounds the bytes and the messages of $u$ that arrive in the remaining $T_1 + T_2$ of $C_i(t)$, with at most $\lceil r_\ell\,(T_1 + T_2 + J_\ell) \rceil$ outputs of topic $\ell$, each with at most $n_\ell$ messages of at most $s_\ell$ bytes;
(b) the held bytes $B_i(t)$, the bytes and the messages of the data bridge of vessel $i$ that the GS has not acknowledged, split by QoS;
(c) the report of step $t-1$: the number of the messages of each admitted unit that the regulator released, the number of the messages of admitted units that the regulator discarded, and whether the table of step $t-1$ reached the vessel.
\end{definition}

The part $(x^{\mathrm{h}}_u, n^{\mathrm{h}}_u)$ of the demand is exact, and the part $(x^{\mathrm{a}}_u, n^{\mathrm{a}}_u)$ is the arrival curve of topic $\ell$ (Section~\ref{subsec:dba}) at the duration $T_1 + T_2$, so $x_u(t)$ exceeds the bytes of $u$ by at most the bytes of the outputs that the arrival curve allows in $T_1 + T_2$.

\begin{definition}[Table of a vessel]\label{def:table}
The table of vessel $i$ in step $t$ is $(t, p_{i,\cdot}(t), k_{i,\cdot}(t))$, where $p_{i,\ell}(t) \in \{0, 1, 2\}$ is the level of topic $\ell$ of vessel $i$ and $k_{i,\ell}(t)$ is the number of the admitted outputs of topic $\ell$ of vessel $i$. Level~2 marks the admitted topics of $\mathcal{N}_i(t)$, level~1 the other admitted topics, and level~0 the withheld topics. The default table admits the units of $\mathcal{R}$ alone.
\end{definition}

\begin{definition}[Vessel regulator]\label{def:regulator}
The vessel regulator of vessel $i$ runs in the MQTT broker of vessel $i$ and holds the messages of step $t$ until $v_i(t)$. From $v_i(t)$ to the end of step $t$ on the clock of vessel $i$, the regulator releases to the data bridge the messages of the admitted units of vessel $i$, first the units of level~2 and then the units of level~1, each in the order $\pi_t$; for a stream topic, up to the demand of its unit; for an output topic, the first $k_{i,\ell}(t)$ complete outputs of step $t$, each up to the demand of its unit. In each tick of period $T_{\mathrm{tick}}$, the regulator releases the next message only while the unsent bytes of the socket of the data bridge plus the bytes that the regulator released in the tick stay below $L_{\mathrm{w}}$. The regulator writes the messages of the topics of $\mathcal{R}$ with QoS~2 and the other messages with QoS~0, each with the index of its step and its sequence number in its topic, and withholds the other messages. At the end of step $t$ on its own clock, the regulator discards the messages of step $t$ that it did not release. A vessel without its table of step $t$ at $v_i(t)$ applies the default table, with the arrival-curve bound over $T_{\mathrm{s}}$ as the demand of each unit, ignores a table of step $t$ that arrives later, and reports the table as not received.
\end{definition}

Let $d_{\mathrm{m}}$ bound the time from the release of a message to its write by the broker, and let $L_{\mathrm{w}} \geq R_{\max}\,(T_{\mathrm{tick}} + d_{\mathrm{m}})/8$, with $R_{\max}$ the largest PHY rate of the radios. After each tick in which the regulator holds an unreleased admitted message, the vessel holds at least $L_{\mathrm{w}}$ bytes of the data bridge that its radio has not transmitted, in the broker, the socket, or the queues below the socket, and the radio transmits at most $R_{\max}(T_{\mathrm{tick}} + d_{\mathrm{m}})/8$ bytes before the next tick. Hence, while the regulator holds an unreleased admitted message, the data bridge of the vessel holds bytes. At the end of step $t$, the socket holds at most $L_{\mathrm{w}}$ plus one message of unsent bytes of step $t$, and $B_i(t+1)$ counts these bytes. The programs of the GS discard each message of step $t$ that arrives after step $t$ and each repeated message, so each message of step $t$ has a firm deadline at the end of step $t$ [kopetz2011real] and reaches the applications of the GS at most once, also when the data bridge reconnects.

\subsection{Prediction of the Capacity}\label{subsec:pred}

The monitor of the AP records, for each frame that the AP sends or receives, its time, its size, its PHY rate, and its station, and records the busy time of the channel with a resolution of $1~\mu$s (Section~\ref{subsec:plat_gs}). The quiet interval $Q(t)$ runs from the end of $W(t-1)$ to the start of phase~2 of step $t$ on the clock of the GS. In $Q(t)$, the vessels send the tables of phase~1 and the bytes that the data bridges hold from step $t-1$.

\begin{definition}[External busy fraction]\label{def:busy}
The external busy fraction of step $t$ is $b(t) = \big(T_{\mathrm{busy}}(Q(t)) - T_{\mathrm{own}}(Q(t))\big)/|Q(t)|$, where $T_{\mathrm{busy}}(Q(t))$ is the busy time of the channel in $Q(t)$ and $T_{\mathrm{own}}(Q(t))$ is the time in which the channel carries, in $Q(t)$, the successful frames that the AP exchanges with the vessels, their MAC acknowledgments, and the beacons of the AP.
\end{definition}

The fraction $b(t)$ counts the frames of the transmitters outside the fleet and the collisions of the fleet in $Q(t)$, and $b(t)$ uses only samples of $Q(t)$, which precedes $W(t)$. With the efficiency $\eta(n)$ of~\eqref{eq:eff}, computed with the EDCA parameters of AC\_BE of the stations, the PHY rate of the DTP, and the mean frame of the data bridges, which is the mean of $T_{\mathrm{succ}}$ over the frames that the topics of the DTP produce at their rates, weighted by the number of frames, the AS predicts the capacity of step $t$ with
\begin{equation}
\hat{\kappa}(t) = \big(1 - \hat{b}(t)\big)\,\eta(M + 1), \qquad \hat{b}(t) = (1 - \lambda)\,\hat{b}(t-1) + \lambda\, b(t),
\label{eq:pred}
\end{equation}
where $\lambda \in (0, 1]$ is the gain of the exponential smoothing (Section~\ref{subsec:aci}), $\hat{b}(t) = \hat{b}(t-1)$ when the monitor provides no $b(t)$, and $M + 1$ counts the $M$ vessels, which release at least the units of $\mathcal{R}$ in every step, and the AP, which sends the TCP acknowledgments. Model~\eqref{eq:eff} assumes stations with one frame size and one rate that hear each other and that offer frames in every slot, and the AP of Fleet-DT sends the acknowledgments of every vessel, the AP uses a shorter AIFS and a smaller contention window than the stations (PR7), the vessels use different rates, and two vessels can be hidden from each other; the scores below measure the error of~\eqref{eq:pred} under these conditions.

The score of step $t$ is $z(t) = \kappa(t)/\hat{\kappa}(t)$, for each step with $e(t) - w(t) \geq T_{\min}$ in which the GS obtains $e(t)$. With the $N^z_t$ scores of the DTA and of the steps before $t$ sorted as $z_{(1)} \leq \dots \leq z_{(N^z_t)}$, the factor $Z_t(\alpha) = z_{(\lceil \alpha N^z_t \rceil)}$ for $\alpha \in (0, 1]$, $Z_t(\alpha) = 1$ for $\alpha > 0$ when $N^z_t = 0$, and $Z_t(\alpha) = 0$ for $\alpha \leq 0$ is nondecreasing in $\alpha$. The DTA provides the initial scores of an operation.

\begin{definition}[Failure indicator and level]\label{def:aci}
The failure indicator $f(t) \in \{0, 1\}$ of step $t$ is $1$ when, for some vessel, the GS accepts within step $t$ fewer messages of an admitted unit than the vessel reports as released, the vessel reports a discarded message of an admitted unit, or a vessel of $\mathcal{M}(t)$ reports that its table of step $t$ did not reach it; or when a vessel of $\mathcal{M}(t)$ sends no report of step $t$ in phase~1 of step $t+1$. Otherwise $f(t) = 0$. With the target $\beta \in (0, 1)$, the step $\gamma \in (0, 1)$, and the initial level $\alpha_1 \in [0, 1]$, the AS computes in phase~2 of step $t+1$ the level
\begin{equation}
\alpha_{t+1} = \min\Big\{1,\ \max\big\{-\gamma(1 - \beta),\ \alpha_t + \gamma\,(\beta - f(t))\big\}\Big\}.
\label{eq:aci}
\end{equation}
\end{definition}

Equation~\eqref{eq:aci} is the update of ACI (Section~\ref{subsec:aci}) with the event ``the GS misses a message of an admitted unit'' in place of the miscoverage event, projected on $[-\gamma(1-\beta), 1]$. A step with $f(t) = 1$ lowers the level by $\gamma(1 - \beta)$, a step with $f(t) = 0$ raises the level by $\gamma\beta$, and the projection bounds the number of steps that the level needs to cross $[0, 1]$ after a change of the channel.

\subsection{Admission Control and Guarantee}\label{subsec:admission}

For a vessel $i \notin \mathcal{M}(t)$, the AS uses as held bytes $\bar{B}_i(t) = \min\{\bar{B}_i(t-1) + x^{\mathrm{adm}}_i(t-1),\ B_{\mathrm{cap}}\}$, where $x^{\mathrm{adm}}_i(t-1)$ holds, for each QoS, the content and the PUBLISH headers of the units that the AS admitted for vessel $i$ in step $t-1$ and, for QoS~2, $h_{\mathrm{A}}$ bytes per message for the PUBREL packet, and $B_{\mathrm{cap}}$ is the send buffer of the data bridge; $\bar{B}_i(t) = B_i(t)$ for $i \in \mathcal{M}(t)$, and $\bar{B}_i(t) = 0$ while the broker of the GS reports the data bridge of vessel $i$ as disconnected. The cost $c^B_i(t)$ of the held bytes applies $a_i$ to $\bar{B}_i(t)$ with the QoS of each part, plus the QoS~2 control packets of the messages of step $t-1$ whose PUBREL has not reached the broker of the GS, which the AS counts as the QoS~2 messages that vessel $i$ released in step $t-1$ minus the QoS~2 messages that the GS accepted minus the QoS~2 messages of $B_i(t)$.

\begin{definition}[Budget and order of the units]\label{def:grant}
The budget of step $t$ is
\begin{equation}
K(t) = Z_t(\alpha_t)\,\hat{\kappa}(t)\,T_W - \sum_{i=1}^{M} c^B_i(t),
\label{eq:grant}
\end{equation}
where the sum runs over the vessels whose data bridge is connected. The order $\pi_t$ sorts the units of step $t$ by the key $\big(\mathbb{1}[\ell \notin \mathcal{R}],\ \mathbb{1}[\ell \notin \mathcal{N}_i(t)],\ \mathrm{rank}_{\sigma_t}(\ell),\ k,\ (i - t) \bmod M\big)$ in increasing lexicographic order.
\end{definition}

The order $\pi_t$ places the units of $\mathcal{R}$ first, then the needed units, then the topics by rank, then the outputs by index, and rotates the vessels by one position per step, so, among the vessels with the same need of topic $\ell$, output $k$ of every vessel precedes output $k+1$ of any vessel.

\begin{algorithm}[t]
\caption{Admission control of phase~2 of step $t$}\label{alg:admission}
\begin{algorithmic}[1]
\Require demands, $B_i(t)$, and reports of the vessels of $\mathcal{M}(t)$; $b(t)$; $\mathcal{N}_i(t)$, $\sigma_t$; $\hat{b}(t-1)$, $\alpha_{t-1}$, $f(t-1)$, the scores, $\bar{B}_i(t-1)$; $\lambda$, $\gamma$, $\beta$, $T_W$
\Ensure $\hat{\kappa}(t)$, $\alpha_t$, $K(t)$, $A(t)$, the tables of the vessels of $\mathcal{M}(t)$
\State $\hat{b}(t)$ and $\hat{\kappa}(t)$ by~\eqref{eq:pred}; $\alpha_t$ by~\eqref{eq:aci}; $K(t)$ by~\eqref{eq:grant}
\State $A \gets$ the units of the topics of $\mathcal{R}$ of the vessels with a connected data bridge, with the arrival-curve bound over $T_{\mathrm{s}}$ as the demand of each vessel outside $\mathcal{M}(t)$
\State $c_A \gets \sum_{u \in A} c_u(t)$
\For{each unit $u = (i, \ell, k) \notin A$ of a vessel $i \in \mathcal{M}(t)$ with a connected data bridge, in the order $\pi_t$}
  \If{$c_A + c_u(t) \leq K(t)$ and ($k = 1$ or $(i, \ell, k-1) \in A$)}\label{alg:test}
    \State $A \gets A \cup \{u\}$; $c_A \gets c_A + c_u(t)$
  \EndIf
\EndFor
\State $p_{i,\ell}(t) \gets 2$ if $(i,\ell,1) \in A$ and $\ell \in \mathcal{N}_i(t)$; $1$ if $(i,\ell,1) \in A$ and $\ell \notin \mathcal{N}_i(t)$; $0$ otherwise
\State $k_{i,\ell}(t) \gets |\{k : (i, \ell, k) \in A\}|$; publish the table of each vessel $i \in \mathcal{M}(t)$ on its topic \texttt{levels}
\end{algorithmic}
\end{algorithm}

Fig.~\ref{fig:alg} shows an example of Algorithm~\ref{alg:admission}. The test of line~\ref{alg:test} admits an output only after the previous output of the same topic and vessel, so the admitted outputs of each topic of each vessel form a prefix, and the count $k_{i,\ell}(t)$ identifies them.

\begin{lemma}\label{lem:greedy}
Algorithm~\ref{alg:admission} returns the set $A(t)$ of Problem~\ref{prob:alloc} for the budget $K(t)$, among the sets whose admitted outputs of each topic of each vessel form a prefix.
\end{lemma}
\begin{proof}
Let $\mathcal{F}$ be the family of the sets of units that hold the units of $\mathcal{R}$, satisfy~\eqref{eq:feas}, and hold a prefix of the outputs of each topic of each vessel. If the units of $\mathcal{R}$ exceed $K(t)$, the test of line~\ref{alg:test} fails for every unit, and $A(t)$ holds the units of $\mathcal{R}$ alone. Otherwise, $\mathcal{F}$ is nonempty, and the key of $\pi_t$ orders the units totally, so the lexicographically largest set $A^* \in \mathcal{F}$ exists and is unique. Let $u_j$ be the $j$-th unit of the loop and $A_j$ the set after iteration $j$, and assume that $A_{j-1}$ holds exactly the units of $A^*$ among $\mathcal{R}$ and $u_1, \dots, u_{j-1}$, which holds for $j = 1$. If $u_j \in A^*$, then $A_{j-1} \cup \{u_j\} \subseteq A^*$; the costs are nonnegative and the outputs of $A^*$ that precede $u_j$ in its topic are in $A_{j-1}$, so the test holds and the algorithm admits $u_j$. If $u_j \notin A^*$, then $A_{j-1} \cup \{u_j\}$ agrees with $A^*$ before $u_j$ and holds $u_j$, so it is lexicographically larger than $A^*$ and does not belong to $\mathcal{F}$, and the test fails. By induction on $j$, the algorithm returns $A^*$.
\end{proof}

Lemma~\ref{lem:greedy} is the lexicographic optimality of the greedy selection over a family that is closed under the removal of the last admitted unit of a topic.

Let the work $X(t)$ of step $t$ be the sum of $T_{\mathrm{succ}}$ over every exchange that the data bridges need, from $w(t)$ on, to deliver the bytes that they hold at $w(t)$ and every message of the units of $A(t)$: the segments, their TCP acknowledgments, the QoS~2 control packets, the retransmissions, and the window updates.

\begin{lemma}\label{lem:delivery}
Under Assumption~\ref{as:clock}, if $X(t) \leq \kappa(t)\,T_W$, the regulators release every message of every unit of $A(t)$, the GS accepts each of them within step $t$, and $B_i(t+1) = 0$ for every vessel $i$.
\end{lemma}
\begin{proof}
Suppose that a data bridge holds bytes at the end of $W(t)$. Then $e(t)$ is the end of $W(t)$ by Definition~\ref{def:capacity}, so the data bridges complete in $W(t)$ exchanges with the airtime $T_{\mathrm{fleet}}([w(t), e(t))) = \kappa(t)\,T_W \geq X(t)$, all of them part of the work $X(t)$, and the exchanges that remain make the work larger than $X(t)$, a contradiction. Hence every data bridge completes its exchanges, including the PUBREL packets of the messages of QoS~2, before the end of $W(t)$. A regulator that holds an unreleased admitted message keeps bytes in its data bridge (Section~\ref{subsec:protocol}), so every regulator releases every admitted message before the end of $W(t)$, which, by Assumption~\ref{as:clock}, precedes the end of step $t$ on the clock of the vessel, and discards no admitted message. By Definition~\ref{def:window}, the broker of the GS delivers each message no later than $d_{\mathrm{rx}}$ after its exchange, before the end of step $t$. The GS acknowledges every byte of the data bridges before the end of $W(t)$, and the regulators release nothing until $v_i(t+1)$, so $B_i(t+1) = 0$.
\end{proof}

Lemma~\ref{lem:delivery} follows from Definition~\ref{def:capacity} and from the pacing of the regulator, and its hypothesis is the operational content: the work satisfies $X(t) \leq \sum_{u \in A(t)} c_u(t) + \sum_{i} c^B_i(t) \leq Z_t(\alpha_t)\,\hat{\kappa}(t)\,T_W$ when the bound of Definition~\ref{def:cost} holds, so the hypothesis holds when $Z_t(\alpha_t)\,\hat{\kappa}(t) \leq \kappa(t)$. The grant charges the held bytes $c^B_i(t)$, so the condition of Lemma~\ref{lem:delivery} involves step $t$ alone, and a step that satisfies it after a shortfall delivers its units and empties the data bridges.

\begin{theorem}\label{thm:aci}
For every realization of the channel, of the demands, and of the predictions, and for every $T \geq 1$,
\begin{equation}
\frac{1}{T}\sum_{t=1}^{T} f(t) \;\leq\; \beta + \frac{\alpha_1}{\gamma T} + \frac{(1 - \beta)\,\big(1 + F_{\mathcal{R}}(T)\big)}{T},
\label{eq:thm}
\end{equation}
where $F_{\mathcal{R}}(T)$ is the number of the steps $t \leq T$ with $\alpha_t \leq 0$ and $f(t) = 1$.
\end{theorem}
\begin{proof}
Let $D(t) = \alpha_{t+1} - \alpha_t - \gamma(\beta - f(t))$ be the correction of the projection in~\eqref{eq:aci}. Summing $\alpha_{t+1} = \alpha_t + \gamma(\beta - f(t)) + D(t)$ over $t = 1, \dots, T$ gives
\[
\frac{1}{T}\sum_{t=1}^{T} f(t) = \beta + \frac{\alpha_1 - \alpha_{T+1} + \sum_{t=1}^{T} D(t)}{\gamma T}.
\]
The projection gives $\alpha_{T+1} \geq -\gamma(1 - \beta)$. The upper limit gives $D(t) \leq 0$. The lower limit gives $D(t) > 0$ only when $\alpha_t + \gamma(\beta - f(t)) < -\gamma(1-\beta)$, which requires $f(t) = 1$ and $\alpha_t < 0$, and then $D(t) = -\gamma(1-\beta) - \alpha_t + \gamma(1-\beta) = -\alpha_t \leq \gamma(1 - \beta)$, because $\alpha_t \geq -\gamma(1-\beta)$. Hence $\sum_{t=1}^{T} D(t) \leq \gamma(1-\beta) F_{\mathcal{R}}(T)$, and $\alpha_1 - \alpha_{T+1} + \sum_t D(t) \leq \alpha_1 + \gamma(1-\beta)(1 + F_{\mathcal{R}}(T))$, which gives~\eqref{eq:thm}.
\end{proof}

The proof follows Gibbs and Cand\`es [gibbs2021adaptive] on the lower side of the level. In ACI, a level below zero yields the whole space as prediction set and forces $\mathrm{err}_t = 0$; in Fleet-DT, a level $\alpha_t \leq 0$ yields $Z_t(\alpha_t) = 0$ and $K(t) \leq 0$, so Algorithm~\ref{alg:admission} admits the units of $\mathcal{R}$ alone, and $f(t)$ can still be $1$. The term $F_{\mathcal{R}}(T)$ counts these steps, and the AS observes $F_{\mathcal{R}}(T)$ in every step.

\begin{corollary}\label{cor:minimal}
Under Assumption~\ref{as:clock}, $F_{\mathcal{R}}(T) \leq H_{\mathcal{R}}(T)$, where $H_{\mathcal{R}}(T)$ counts the steps $t \leq T$ in which the AS admits the units of $\mathcal{R}$ alone and $X(t) > \kappa(t)\,T_W$, or a table or a report of step $t$ is lost.
\end{corollary}
\begin{proof}
In a step counted by $F_{\mathcal{R}}(T)$, $A(t)$ holds the units of $\mathcal{R}$ alone and $f(t) = 1$. If no table and no report of step $t$ is lost, $f(t) = 1$ means that the GS misses a released message or a regulator discards an admitted message, and both cases contradict the conclusion of Lemma~\ref{lem:delivery}, so $X(t) > \kappa(t)\,T_W$.
\end{proof}

\begin{corollary}\label{cor:state}
Let the topic \texttt{state} belong to $\mathcal{R}$, with one output per step whose release stays more than $J_{\mathrm{state}}$ away from $v_i(t)$. In every step $t$ that satisfies the hypothesis of Lemma~\ref{lem:delivery}, the GS accepts within step $t$, exactly once, the state of step $t$ of every vessel of $\mathcal{M}(t)$, at most $T_{\mathrm{s}} + T_3 + \varepsilon$ after its arrival at the regulator, and over a sequence of such steps the age of information of the state of each vessel at the GS, counted from the arrival at the regulator, stays below $2T_{\mathrm{s}} + T_3 + \varepsilon$. The fraction of the steps $t \leq T$ in which the GS misses the state of a vessel of $\mathcal{M}(t)$ is at most the right side of~\eqref{eq:thm}.
\end{corollary}
\begin{proof}
Algorithm~\ref{alg:admission} admits the units of $\mathcal{R}$ of every vessel in every step, the release condition places one state in each collection window, Lemma~\ref{lem:delivery} gives the acceptance within step $t$, the programs of the GS discard the repeated messages, and Section~\ref{subsec:window} gives the delay. Just before the acceptance of the state of step $t+1$, no later than $t_0 + (t+1)\,T_{\mathrm{s}}$, the newest accepted state is the state of step $t$, which arrived at its regulator no earlier than $v_i(t-1) \geq t_0 + (t-2)\,T_{\mathrm{s}} + T_1 + T_2 - \varepsilon$ on the clock of the GS, which gives the age bound; the delay from the source to the regulator adds to both bounds. A step in which the GS misses a state has $f(t) = 1$, and Theorem~\ref{thm:aci} bounds the fraction of these steps.
\end{proof}

\begin{definition}[Slotted data window]\label{def:tdma}
In the slotted variant of Fleet-DT, the AS divides $W(t)$ into a shared slot followed by one slot per vessel of $\mathcal{M}(t)$ in the rotated order of $\pi_t$. In the shared slot, every vessel sends the units of $\mathcal{R}$ and its held bytes, and the shared slot lasts the cost of these units and bytes divided by $Z_t(\alpha_t)\hat{\kappa}(t)$. The slot of vessel $i$ lasts the cost of its admitted units outside $\mathcal{R}$ divided by $Z_t(\alpha_t)\hat{\kappa}_2(t)$, where $\hat{\kappa}_2(t)$ is~\eqref{eq:pred} with $\eta(2)$, plus a guard of $2\varepsilon$ plus the airtime of the largest exchange. The vessel regulator releases the messages of the admitted units of vessel $i$ outside $\mathcal{R}$ only in the slot of vessel $i$, and the radio of vessel $i$ holds at most one frame of the data bridge in its queue.
\end{definition}

In the slotted variant, Algorithm~\ref{alg:admission} admits a unit when the duration of the shared slot plus the durations of the slots of the vessels, with the unit, stays within $T_W$; this duration grows with every admitted unit, so Lemma~\ref{lem:greedy} holds for the slotted variant. In each slot of a vessel, one vessel and the AP contend. The capacity $\kappa_j(t)$ of slot $j$ follows Definition~\ref{def:capacity} with the slot in place of $W(t)$, and Lemma~\ref{lem:delivery} applies to each slot with $\kappa_j(t)$, the duration of the slot, and the work of the slot. Section~\ref{sec:method} compares the slotted variant with Fleet-DT.

\subsection{Digital Twin of the Network}\label{subsec:dts}

Fleet-DT twins the network of Section~\ref{subsec:network} with the DTP, the DTI, and the DTA of Grieves and Vickers and with a DTE that holds the application space (AS) and the virtual environment (VE).

\begin{definition}[Physical unit]\label{def:pu}
The physical unit (PU) is the network that carries the messages of one operation: the channel, the radios, the MQTT brokers, and the data bridge and the control bridge of each vessel.
\end{definition}

\begin{definition}[Graph of the network]\label{def:graph}
The graph of the network in step $t$ is the attributed directed multigraph $G(t) = (V(t), E(t))$ with the vertices $V(t) = \{0, 1, \dots, M\}$, where vertex~$0$ is the GS, and the edges $E(t) = \{(i, \ell)\}$ from vertex $i$ to vertex~$0$ for $i \in \{1, \dots, M\}$ and $\ell \in \{1, \dots, L\}$. The edge $(i, \ell)$ carries the demand of the units of topic $\ell$ of vessel $i$, $p_{i,\ell}(t)$, $k_{i,\ell}(t)$, the released and the discarded messages, and the bytes $d_{i,\ell}(t)$ that the GS accepts within step $t$. Vertex $i$ carries $R_i(t)$, $R_{0,i}(t)$, $\bar{B}_i(t)$, and the membership of $i$ in $\mathcal{M}(t)$. Vertex~$0$ carries $b(t)$, $\hat{b}(t)$, $\hat{\kappa}(t)$, $\kappa(t)$, $e(t)$, $z(t)$, $\alpha_t$, $K(t)$, and $f(t)$.
\end{definition}

\begin{definition}[DTP]\label{def:dtp}
The DTP is the model of the architecture of the network without the values of any step: the roles of the hosts; the PHY and EDCA parameters of the channel; the two bridges of each vessel with their user priorities; the uplink topics with $r_\ell$, $J_\ell$, $n_\ell$, $s_\ell$, the kind, and the QoS; the required set $\mathcal{R}$; the downlink topic \texttt{goals} with its rate and size bound; the cost function $a_i$ (Section~\ref{subsec:costplat}); $T_{\mathrm{s}}$, $T_1$, $T_2$, $\varepsilon$, $d_{\mathrm{rx}}$, $d_{\mathrm{m}}$, $T_{\mathrm{tick}}$, $L_{\mathrm{w}}$, $B_{\mathrm{cap}}$, $T_{\min}$; $\beta$; the initial order and needed sets; and the types of the vertices, the edges, and the attributes of $G(t)$.
\end{definition}

\begin{definition}[DTI]\label{def:dti}
The DTI is the instance of the DTP linked to the PU, and the DTI holds the state of the network in each step $t$ as the graph $G(t)$ with the values of its attributes. The DTI program of each vessel computes and publishes the values of Definition~\ref{def:demand} in phase~1, and the GS holds the other values of the DTI.
\end{definition}

\begin{definition}[DTA]\label{def:dta}
The DTA is the aggregation of the DTIs of the past operations and of the emulated runs, one sequence of graphs $G(t)$ per operation and per run.
\end{definition}

\begin{definition}[DTE]\label{def:dte}
The DTE is the set of the programs of the GS that operate on the DTI and on the DTA. The AS holds the models of the network: the efficiency model~\eqref{eq:eff} and the prediction~\eqref{eq:pred}, the level~\eqref{eq:aci} and the budget~\eqref{eq:grant}, Algorithm~\ref{alg:admission}, the fidelity of the DT (Definition~\ref{def:fidelity}), and the what-if evaluation (Definition~\ref{def:whatif}). The VE draws $G(t)$ for the operator in each step.
\end{definition}

The DTI mirrors the values that the PU measures (the demands, $\bar{B}_i$, the reports, the PHY rates, $b$, $\kappa$, $e$, and $d_{i,\ell}$) and holds the values that the AS predicts or decides ($\hat{\kappa}$, $\alpha_t$, $K$, $A$, and the tables). The PU and the DTI exchange data in both directions in every step (Table~\ref{tab:exchanges}). In the terms of Zhou et al.~[zhou2026ndt], the DTI and the DTA hold the data of the NDT, the AS holds the models, the control bridges and the monitor of the AP realize the mapping, and the broker of the GS realizes the interface to the operation control software, which provides $\mathcal{N}_i(t)$ and $\sigma_t$ and stays outside the DT of the network.

\subsection{Fidelity and What-If Evaluation}\label{subsec:fidelity}

\begin{definition}[Fidelity of the DT]\label{def:fidelity}
Let $S$ be a set of steps of the DTA. The fidelity of the DT on $S$ comprises: (i) the failure fraction $\bar{f} = \frac{1}{|S|}\sum_{t\in S} f(t)$ against $\beta$; (ii) the coverage of the bound, the fraction of the steps $t \in S$ with $\kappa(t) \geq Z_t(\alpha_t)\,\hat{\kappa}(t)$; (iii) the mean absolute error of $\hat{\kappa}(t)$ against $\kappa(t)$; (iv) the ratio of $d_{i,\ell}(t)$ to the released bytes of each edge with $p_{i,\ell}(t) \geq 1$; (v) the utilization $\nu = \sum_{t \in S}\sum_{u \in A(t)} c_u(t) \big/ \sum_{t\in S} \kappa(t)\,T_W$; and (vi) the age of information of the state of each vessel at the GS [kaul2012real].
\end{definition}

\begin{definition}[What-if evaluation]\label{def:whatif}
For a configuration of the network, given by $M$, the PHY rates, the topics, $T_{\mathrm{s}}$, and $\beta$, the AS predicts $\bar{f}$ and $\nu$ by replaying the scores and the external busy fractions of the DTA through~\eqref{eq:pred}, \eqref{eq:aci}, and Algorithm~\ref{alg:admission} with the efficiency $\eta$ of the configuration, under the assumption that the scores and the external busy fractions of the DTA hold for the configuration.
\end{definition}

Fig.~\ref{fig:loop} summarizes Section~\ref{sec:dt} as a closed loop. In the inner loop, once per step, the DTI receives the demands, the held bytes, the reports, and the records of the monitor of the AP, and the AS computes $\hat{\kappa}(t)$, $\alpha_t$, $K(t)$, and $A(t)$ and sends the tables to the regulators. In the outer loop, between operations, the AS sets $\lambda$, $\gamma$, and the initial scores from the DTA, evaluates configurations with Definition~\ref{def:whatif}, and measures the fidelity of the DT with Definition~\ref{def:fidelity}; $\gamma$ stays fixed within an operation, as Theorem~\ref{thm:aci} requires.

\begin{figure}[t]\centering\footnotesize\resizebox{\columnwidth}{!}{\input{fig_admission}}\caption{Example of Algorithm~\ref{alg:admission} with $M = 2$, $\mathcal{R} = \{\texttt{state}\}$, $\mathcal{N}_1(t) = \{\texttt{state}, \texttt{imu}\}$, $\mathcal{N}_2(t) = \{\texttt{state}, \texttt{imu}, \texttt{video}\}$, $\sigma_t = (\texttt{state}, \texttt{imu}, \texttt{gps}, \texttt{video}, \texttt{baro})$, and $K(t) = 40$~ms.}\label{fig:alg}\end{figure}
\begin{figure*}[t]\centering\resizebox{0.85\textwidth}{!}{\input{fig_loop}}\caption{Closed loop of the DBA and the constructs of the DT of the network.}\label{fig:loop}\end{figure*}
\section{Platform Requirements}\label{sec:platform}

This section maps the components of Section~\ref{sec:dt} to the programs and the devices of a platform of Fleet-DT and states the minimum requirements that these programs and devices meet. Fig.~\ref{fig:deploy} shows the programs and the devices as a UML deployment diagram, and Table~\ref{tab:req} lists the requirements with the element of Section~\ref{sec:dt} that each requirement serves.

\begin{table}[t]
\caption{Minimum requirements of a platform of Fleet-DT.}\label{tab:req}
\centering\footnotesize
\begin{tabular}{@{}lp{5.4cm}l@{}}
\toprule
& Requirement & Serves \\
\midrule
PR1 & The MQTT broker of each vessel runs a plugin that receives each published message, republishes a message, wakes at the period $T_{\mathrm{tick}}$ and at given instants, and reads the unsent and the unacknowledged bytes of the socket of the data bridge. & Defs.~\ref{def:demand}, \ref{def:regulator} \\
PR2 & The DTP declares, for each topic, $r_\ell$, $J_\ell$, $n_\ell$, $s_\ell$, the kind, and the QoS. & Defs.~\ref{def:topic}, \ref{def:demand} \\
PR3 & The DTI program of each vessel publishes the values of Definition~\ref{def:demand} before the end of phase~1. & Def.~\ref{def:demand} \\
PR4 & The clock of each vessel and of the radio of the GS follows the clock of the GS with an offset of at most $\varepsilon$, with slew and without backward steps. & Assumption~\ref{as:clock} \\
PR5 & The broker of the GS opens, for each vessel, one data bridge and one control bridge over their own TCP connections, and the hosts give UP~6 to the packets of the control bridge. & Section~\ref{subsec:window} \\
PR6 & The radio of the GS is the AP, and its driver reports the busy time of the channel with a resolution of $1~\mu$s and the time, the size, the PHY rate, and the station of each frame that the AP sends or receives. & Defs.~\ref{def:capacity}, \ref{def:busy} \\
PR7 & The radios use EDCA and OFDM rates, send each IP packet in one frame without aggregation, and the AP uses for its own AC\_BE a shorter AIFS and a smaller contention window than the stations. & Def.~\ref{def:cost}, \eqref{eq:pred} \\
PR8 & The TCP connections of the bridges send without the Nagle algorithm on both ends, keep the congestion window across phases~1 and~2, and use a minimum retransmission timeout of at most $T_3/4$. & Def.~\ref{def:cost}, Def.~\ref{def:capacity} \\
PR9 & The MQTT brokers write each forwarded message to the socket of the bridge without a limit on the messages in flight. & Def.~\ref{def:regulator} \\
\bottomrule
\end{tabular}
\end{table}

\subsection{Vessel}\label{subsec:plat_vessel}
Each vessel carries a computer and a radio that connect through a network interface. The computer runs the vessel MQTT broker, the sources of the uplink topics, and the DTI program, each with its own MQTT client. The sources publish on the local namespace \texttt{boat/b$i$/raw}, which no bridge forwards, and the vessel regulator, a plugin of the vessel MQTT broker (PR1), republishes each released message on the namespace \texttt{boat/b$i$/tx}, which the data bridge forwards. The regulator finds the socket of the data bridge by the local port of the socket, because the data bridge connects to a dedicated listener port of the vessel MQTT broker. In each tick (the event \texttt{MOSQ\_EVT\_TICK} of the plugin interface, with the next wake-up at most $T_{\mathrm{tick}}$ later), the regulator reads with \texttt{SIOCOUTQNSD} the unsent bytes of this socket and releases messages while these bytes plus the bytes released in the tick stay below $L_{\mathrm{w}}$ (Definition~\ref{def:regulator}). In phase~1, the regulator reads with \texttt{SIOCOUTQ} the bytes of this socket that the GS has not acknowledged and, from the offsets of the messages that it released, the number of these messages and their QoS. A source of an output topic splits each output into messages of at most $s_\ell$ bytes, and the content of each message starts with a header of 8 bytes that holds the index of the output, the index of the message in the output, and the number of messages of the output. The operating system runs the vessel MQTT broker and the DTI program with a real-time scheduling priority.

\subsection{Ground Station}\label{subsec:plat_gs}
The GS consists of a computer and a radio that connect through a network interface. The computer runs the GS MQTT broker, the Fleet-DT program, which holds the DTP, the part of the DTI at the GS, and the DTA and runs the AS and the VE, and the operation control software, which publishes $\mathcal{N}_i(t)$ and $\sigma_t$ in phase~1 on the topic \texttt{fleet/order}. The Fleet-DT program publishes the table of each vessel of $\mathcal{M}(t)$ with QoS~1 on the topic \texttt{levels} of the vessel. The computer of the GS is the time server of the vessels (PR4). The monitor of the AP runs on the radio of the GS, which exchanges every data frame of the fleet, and publishes, at the end of $Q(t)$ and at the end of $W(t)$, the busy time of the channel and the records of the frames of the interval (PR6). Each record carries the access category of the frame, so the AS separates the frames of the data bridges (AC\_BE) from the frames of the control bridges (AC\_VO). From these records, the AS computes $b(t)$, $e(t)$, $\kappa(t)$, and the PHY rates $R_i(t)$ and $R_{0,i}(t)$ of the last frames of each vessel; an exchange that starts in an interval and ends after it counts in proportion to its overlap with the interval. The AS receives the frames of the data bridges in a client and a thread of their own, so the load of the AS stays outside $d_{\mathrm{rx}}$.

\subsection{Network and Airtime Cost}\label{subsec:costplat}
The data bridge subscribes with QoS~2, so each message keeps the QoS of its topic (Definition~\ref{def:required}), and the control bridge subscribes with QoS~1. Both bridges connect with a clean session, so a reconnection resends no message of an earlier step. The vessel regulator writes in the first 6 bytes of the payload of each released message the index of its step (u32) and its sequence number in its topic and step (u16).

With the OFDM PHY of IEEE~802.11, the frame of an IP packet of $s$ bytes occupies $T_{\mathrm{PRE}} + T_{\mathrm{SIG}} + T_{\mathrm{SYM}}\lceil (22 + 8(h_{\mathrm{MAC}} + s))/N_{\mathrm{DBPS}} \rceil + T_{\mathrm{SE}}$, where $h_{\mathrm{MAC}}$ holds the MAC header with the three-address format of an infrastructure BSS and the QoS control field, LLC, and FCS, $N_{\mathrm{DBPS}}$ is the number of data bits per symbol of the rate $R = N_{\mathrm{DBPS}}/T_{\mathrm{SYM}}$, and $T_{\mathrm{SE}}$ is the signal extension of ERP-OFDM [ieee80211std2020]. Bounding the ceiling by its argument plus one gives
\begin{equation}
T_{\mathrm{succ}}(s, R) \leq T_{\mathrm{oh}}(R) + 8s/R,
\label{eq:tsucc}
\end{equation}
where $T_{\mathrm{oh}}(R) = T_{\mathrm{PRE}} + T_{\mathrm{SIG}} + T_{\mathrm{SYM}} + (22 + 8h_{\mathrm{MAC}})/R + T_{\mathrm{SE}} + \mathrm{SIFS} + T_{\mathrm{ACK}}(R) + \mathrm{AIFS}$, $T_{\mathrm{ACK}}(R)$ is the duration of the MAC acknowledgment at the control response rate of $R$, and AIFS is the AIFS of AC\_BE of the stations, which is not shorter than the AIFS of the AP (PR7).

Let $h_{\mathrm{T}}$ be the largest TCP and IP header of one segment, with the options of the connection, $h_{\mathrm{P}}$ the largest PUBLISH header among the uplink topics plus the 6 bytes of step index and sequence number, and $h_{\mathrm{A}} = 4$ the bytes of one MQTT QoS~2 control packet (PUBREC, PUBREL, or PUBCOMP). The regulator writes each message with one call, and TCP sends without the Nagle algorithm (PR8), so $n$ messages of topic $\ell$ with $x$ bytes of content occupy at most $N_{\mathrm{seg}}(x, n) = \min\{n + (x + n h_{\mathrm{P}})/\mathrm{MSS},\ n\,\lceil (s_\ell + h_{\mathrm{P}})/\mathrm{MSS} \rceil\}$ segments, and the GS answers each segment with at most one TCP acknowledgment. With~\eqref{eq:tsucc}, the airtime cost of Definition~\ref{def:cost} is
\begin{equation}
\begin{aligned}
a_i(x, n, q) ={}& N_{\mathrm{seg}}(x, n)\,\big[T_{\mathrm{oh}}(R_i) + 8h_{\mathrm{T}}/R_i + T_{\mathrm{oh}}(R_{0,i}) + 8h_{\mathrm{T}}/R_{0,i}\big] + 8\,(x + n h_{\mathrm{P}})/R_i \\
&+ \mathbb{1}[q = 2]\; n\,\big[3\,T_{\mathrm{oh}}(R_{0,i}) + 3\,T_{\mathrm{oh}}(R_i) + 8(3h_{\mathrm{T}} + 2h_{\mathrm{A}})/R_{0,i} + 8(3h_{\mathrm{T}} + h_{\mathrm{A}})/R_i\big],
\end{aligned}
\label{eq:cost}
\end{equation}
where the second line holds the PUBREC and the PUBCOMP of the GS, the PUBREL of the vessel, and the TCP acknowledgment of each of them. A rate in $W(t)$ below $R_i(t)$ or $R_{0,i}(t)$, a TCP window update, or a retransmission adds work outside~\eqref{eq:cost}, and the frames of the clock synchronization, the MQTT keepalive packets, ARP, the beacons, and the goals in $W(t)$ lower $\kappa(t)$; $f(t)$ and the scores record both effects. The hosts set UP~6 on the packets of the control bridge with the socket priority 262 (256 plus the UP), which the Linux IEEE~802.11 stack maps to the UP of the frame (PR5).

\subsection{Bidirectional Exchange and Synchronization}\label{subsec:exchange}
The PU and the DTI exchange data in both directions in every step without an action of the operator (Table~\ref{tab:exchanges}). In content, every datum carries the index $t$ of its step, the AS assigns each datum to the graph $G(t)$ of its step, and the programs of the GS discard each datum whose index differs from the current step, except a demand of step $t$ that arrives up to $\varepsilon$ before step $t$, and each repeated datum. In time, each host starts each phase on its own clock, which follows the clock of the GS (PR4), and the DTI lags the PU by at most one step. A lost or late datum leaves the exchange in the state of the last column of Table~\ref{tab:exchanges}, and the next step starts the exchange from this state.

\begin{table}[t]
\caption{Exchanges between the PU and the DTI in step $t$.}\label{tab:exchanges}
\centering\footnotesize
\begin{tabular}{@{}p{1.7cm}p{2.0cm}p{1.6cm}p{2.1cm}@{}}
\toprule
Datum & From, to, topic & Time & If missing \\
\midrule
Demand, $B_i(t)$, report of $t-1$ & DTI program of $i$ to AS, \texttt{demand} & phase 1 & $i \notin \mathcal{M}(t)$, $\bar{B}_i(t)$ from step $t-1$, $f(t-1) = 1$ \\
Busy time, frames of $Q(t)$ & monitor to AS & start of phase 2 & $\hat{b}(t) = \hat{b}(t-1)$ \\
Frames of $W(t)$ & monitor to AS & end of step $t$ & no score of step $t$ \\
Accepted bytes $y_{i,\ell}(t)$ & data bridge of $i$ to AS & $W(t)$ & no bytes \\
Table of $i$ & AS to regulator of $i$, \texttt{levels} & phase 2 & default table \\
\bottomrule
\end{tabular}
\end{table}

\begin{figure*}[t]\centering\resizebox{0.8\textwidth}{!}{\input{fig_deploy}}\caption{Programs and devices of a platform of Fleet-DT. The data bridge (magenta) and the control bridge (red) of each vessel cross the radios and the IEEE~802.11 channel; the labels PR1 to PR9 refer to Table~\ref{tab:req}.}\label{fig:deploy}\end{figure*}
\section{Experimental Methodology}\label{sec:method}

The experiments run the software of Fleet-DT without changes on emulated hosts over an emulated IEEE~802.11 channel, so the emulator records the ground truth of every frame and every clock. This section states the hypotheses (Section~\ref{subsec:hyp}), the testbed and its validation (Section~\ref{subsec:testbed}), the scenarios and the methods (Section~\ref{subsec:scen}), the metrics (Section~\ref{subsec:metrics}), and the procedure and the statistics (Section~\ref{subsec:proc}). A tagged commit of the repository holds the hypotheses, the thresholds, the seeds, and the analysis scripts before the first evaluation run (Section~\ref{subsec:proc}).

\subsection{Research Questions and Hypotheses}\label{subsec:hyp}

Table~\ref{tab:hyp} maps each hypothesis to the element of the model that it tests, to the experiments of Table~\ref{tab:exp}, and to its threshold.

\begin{table}[t]
\caption{Hypotheses, the element of the model that each hypothesis tests, and the threshold of each hypothesis.}\label{tab:hyp}
\centering\footnotesize
\begin{tabular}{@{}lp{1.3cm}p{4.9cm}l@{}}
\toprule
& Element & Hypothesis and threshold & Exp. \\
\midrule
H1 & Assumption~\ref{as:clock} & After the convergence of the clocks, $|\theta_i(t)| \leq \varepsilon$ in every probe sample of every vessel, and the demands miss phase~1 in at most $10^{-3}$ of the pairs $(t, i)$ of the steps without external interference, pooled over the runs of each scenario. & all \\
H2 & Defs.~\ref{def:capacity}, \ref{def:busy}; Lemma~\ref{lem:delivery} & The values of $\kappa(t)$ and $b(t)$ that the AS computes from the monitor differ from the ground truth by at most $0.01$ at the 99th percentile; every step with $X(t) \leq \kappa(t)\,T_W$ in the ground truth has $f(t) = 0$ unless a table or a report is lost. & E1 \\
H3 (RQ.II) & Eq.~\eqref{eq:pred} & On the evaluation runs, the prediction~\eqref{eq:pred} has a lower mean absolute error of $\kappa(t)$ than persistence ($\hat{\kappa}(t) = \kappa(t-1)$), exponential smoothing of $\kappa$, and the mean of the calibration runs, and, with ACI at the same $\beta$, a larger utilization $\nu$. & E2 \\
H4 (RQ.I) & Theorem~\ref{thm:aci} & Inequality~\eqref{eq:thm} holds at every prefix of every run, and $\bar{f} \leq \beta + 0.02$ in every run with $F_{\mathcal{R}}(T) = 0$. & E2, E3 \\
H5 (RQ.III) & Algorithm~\ref{alg:admission} & For $\rho \geq 1.5$, Fleet-DT delivers within their steps a larger fraction of the bytes of the needed units, and misses fewer states, than B1, B2, B3, B5, and the overlay TDMA. & E3 \\
H6 & Def.~\ref{def:whatif} & The what-if evaluation, from the DTA of E1 ($M = 4$, $T_{\mathrm{s}} = 125$~ms), predicts $\bar{f}$ within $0.02$ and $\nu$ within $0.05$ for the configurations of E4. & E4 \\
H7 & Section~\ref{subsec:window} & The delay and the age of information of the states stay within the bounds of Corollary~\ref{cor:state} in every step that meets the hypothesis of Lemma~\ref{lem:delivery}. & E4 \\
\midrule
H8 (RQ.III) & Algorithm~\ref{alg:admission} & For $\rho \geq 1.5$, Fleet-DT delivers within their steps a larger fraction of the bytes of the needed units, and misses fewer states, than B6. & E3, E5 \\
H9 (RQ.III) & Algorithm~\ref{alg:admission}, Theorem~\ref{thm:aci} & With $M = 8$ and the sources of E3 on every vessel, Fleet-DT delivers within their steps a larger fraction of the bytes of the needed units, and misses fewer states, than B1, B3, B6, and TDMA, and $\bar{f} \leq \beta + 0.02$ in every run with $F_{\mathcal{R}}(T) = 0$. & E6b \\
\bottomrule
\end{tabular}
\end{table}

\subsection{Testbed}\label{subsec:testbed}

Each vessel and the GS run in one KVM virtual machine with Debian~13, one virtual CPU, and their own clock, and each machine runs the programs of Section~\ref{sec:platform} without changes: Mosquitto~2.1.2 with the patches of the repository, the vessel regulator, the sources of the topics, chrony, and, at the GS, the AS and the operation control software. The data interface of each machine reaches the emulator over a UDP pair. The emulator runs ns-3.48 in real time [nsnam2020wifi]: one AP for the GS and one station per vessel, IEEE~802.11g ERP-OFDM with EDCA on a spectrum channel, and a relay that carries the Ethernet frames of each machine as the IP packets of its Wi-Fi node, which is the role of the radio of Section~\ref{sec:platform}. The relay maps the DSCP value EF to UP~6, and the AP uses for its own AC\_BE the AIFS and the contention window of Table~\ref{tab:param} (PR7). The emulator acts as the monitor of the AP (PR6) and publishes the busy time and the records of the frames, with group-addressed frames in $T_{\mathrm{own}}$, to the broker of the GS every 5~ms and 2.5~ms after the end of $Q(t)$ and of $W(t)$. The hosts mark the packets of the clock synchronization with DSCP EF (AC\_VO). The emulator also writes the ground truth: every transmission and reception of every node, with its time, duration, size, rate, and outcome. A probe reads the clock of each machine over vsock every 100~ms and records its offset to the clock of the host.

Table~\ref{tab:valid} gives the validation of the testbed. The saturation throughput of the emulator for $n = 1, \dots, 16$ stations agrees with the reference of ns-3 within 0.93\% and with~\eqref{eq:eff} within 1.6\% under the EDCA parameters of Table~\ref{tab:param}. The values of $b(t)$ from the monitor agree with the ground truth within 0.0055 at the 99th percentile and the values of $\kappa(t)$ within $1.2 \times 10^{-4}$. The lag of the real-time scheduler stays below 0.3~ms at the 99th percentile in S1 and S3 to S5 and below 0.6~ms in S2 for $M \leq 4$; the analysis discards the steps in which the lag exceeds 1~ms and reports their number, and a run with more than 1\% of such steps repeats once. With $M = 8$ or $T_{\mathrm{s}} = 250$~ms, more than 1\% of the steps exceed 1~ms of lag on the host of the experiments, so E4 holds $M = 2$ and $T_{\mathrm{s}} = 62.5$~ms, the latter with the YansWifiPhy model of ns-3, whose $\kappa(t)$ and $b(t)$ differ from the spectrum model by a Kolmogorov--Smirnov distance below 0.034 in S2, and Section~\ref{sec:results} reports the what-if predictions for $M = 8$ and $T_{\mathrm{s}} = 250$~ms without a measured counterpart.

Fleets of more than four vessels run in a discrete-event simulation without real time. The simulator reuses the code of the emulator for the channel, the Wi-Fi nodes, the monitor, and the ground truth; links the C core of the vessel regulator into ns-3 with the loop timing of the Mosquitto broker; runs the AS of Section~\ref{sec:platform} at each phase~2 with the events of the GS and the monitor records; and replays the release instants and sizes of the sources of the emulation runs of the same seed. ns-3 replaces the Linux hosts: TCP of ns-3 with CUBIC, SACK, timestamps, a minimum RTO of 200~ms, and one ACK per segment, MQTT bridges with the QoS exchanges of Mosquitto and the priorities of Table~\ref{tab:param}, clocks without offset, and fixed delays for the hops inside a host. The simulator writes the files of an emulation run, so the analysis of Section~\ref{subsec:proc} applies without changes. The simulated saturation capacity $\kappa_{\mathrm{sat}}$ (B1 at saturation, two seeds) exceeds the emulated one by 0.039 to 0.057 for $M \in \{4, 8\}$, because the simulated frames collide less (in S2 with $M = 4$, 6\% of the TCP ACKs of the AP are retransmissions against 15\% in the emulation). A validation experiment (E6a) runs the E3 and E5 run files of S1 and S2 at $\rho = 1.5$ in the simulator and compares, per method and scenario, the medians of $D_{\mathcal{N}}$, of the state miss ratio, and of $\kappa(t)$ (thresholds 0.05, 0.01, and 0.05), of $\bar{f}$ and $\nu$ for Fleet-DT and B6 (0.02 and 0.05), and the sign of the paired difference between Fleet-DT and each baseline with the emulation; together with $|\Delta\kappa_{\mathrm{sat}}| \leq 0.03$, these checks gate H9, and the $\kappa_{\mathrm{sat}}$ check fails, so Section~\ref{sec:results} reports E6b without a test of H9. In E6b and E6d, every vessel keeps the sources of the runs with $M = 4$ at $\rho = 1.5$, so the offered load grows with $M$, and the phases keep the values of Table~\ref{tab:param}. Pilot runs with a seed outside the evaluation set show that the control bridges in AC\_VO lose most of their frames to collisions with 16 or more vessels, so E6b covers $M = 8$ and E6d describes $M \in \{16, 32\}$.

\subsection{Scenarios and Methods}\label{subsec:scen}

\paragraph*{Sources} Each vessel runs the topics of Table~\ref{tab:topics}: six telemetry streams and one video topic of the kind output, whose output sizes replay, from a start that the seed sets, the JPEG sizes of 4500 stereo frames of the sequence \texttt{pohang00} of the Pohang Canal Dataset, a stereo recording of coastal water and a port from a vessel [chung2023pohang], with each eye cropped to the aspect ratio 672:376, resized to $336 \times 188$ pixels, and encoded side by side in JPEG at the quality that sets $\rho$. The offered-load ratio $\rho$ divides the airtime cost~\eqref{eq:cost} of all units that the fleet offers per step by $\kappa\,T_W$ measured with a saturated fleet in the same scenario; the sources set $\rho$ through the rate and the quality of the video. The operation control software alternates every 10~s between a survey mode, with the video of one designated vessel in its $\mathcal{N}_i(t)$, and a transit mode, with no video in any $\mathcal{N}_i(t)$; $\mathcal{R} = \{\texttt{state}\}$.

\paragraph*{Channel scenarios} The channel of S1 to S4 lies in the 2.4~GHz band, whose free-space path loss is 6.6~dB lower than at 5.18~GHz, and carries every frame at the ERP-OFDM rate of 24~Mb/s, the highest of the three OFDM rates that IEEE~802.11 makes mandatory [ieee80211std2020], so the methods of one cell differ in the allocation and not in the rate. S1 (clean): every node hears every node. S2 (external network): S1 plus a second BSS on the same channel whose station sends UDP bursts with an on/off schedule (CSMA). S3 (jammer): S1 plus a transmitter without carrier sense with a duty-cycle schedule. S4 (hidden vessels): a loss matrix in which two pairs of vessels do not hear each other and every vessel hears the AP. S5 (mobility): the vessels follow the trajectories of a survey mission over water at 101 to 888~m from the AP, with the antenna of the AP at 6~m and the antennas of the vessels at 1.5~m. The maritime channel near the sea surface has few scatterers, and its models combine the direct ray with the ray that the sea surface reflects [wang2018maritime]; S5 uses the two-ray ground model at 2.412~GHz, whose crossover distance of 910~m exceeds the range of the mission, so the mean path loss of S5 follows free space. Measurements of WLAN links of USVs over seawater fit the two-ray model with 2 to 3~dB of loss above the loss over land at 2.4~GHz and antennas at 2~m [yamamoto2019rssi], a loss that S5 omits. The small-scale fading of S5 follows the Nakagami distribution with $m = 5$, which matches a Rician factor $K$ of 9.3~dB under $m = (K+1)^2/(2K+1)$, below the Rician factors of 17.6 to 18.8~dB that Sun et al.\ measure on land-to-ship links at 5.8~GHz and 3 to 12~km [sun2026landship], so the fading of S5 is deeper than the measured fading. Minstrel adapts the rate of S5 over the OFDM rates of IEEE~802.11a, because the Minstrel model of ns-3 also probes the DSSS rates of IEEE~802.11g, which the radios of PR7 do not use.

\paragraph*{Methods} Fleet-DT (Algorithm~\ref{alg:admission} with ACI), Fleet-DT with the slotted data window (Definition~\ref{def:tdma}), and seven baselines: B1, one data bridge per vessel without regulation; B2, one data bridge per class with the DSCP values EF, AF41, and default, which EDCA maps to AC\_VO, AC\_VI, and AC\_BE; B3, strict priority in the order $\sigma_t$ with the discard at the end of the step, without admission; B4, Algorithm~\ref{alg:admission} with a constant factor $Z$ in place of $Z_t(\alpha_t)$, per scenario the $\beta$-quantile of the scores of the calibration runs divided by the factor $c_S$ of the calibration (Section~\ref{subsec:proc}); B5, Fleet-DT without the policing and the discard of the regulator; TDMA, an overlay TDMA with equal slots per vessel in $W(t)$ and no admission [rao2005overlay, djukic2009soft]; and B6, the interleaved polling with limited service of IPACT [kramer2002ipact] in $W(t)$: the units of $\mathcal{R}$ and the held bytes take a shared slot as in Definition~\ref{def:tdma}, each vessel then receives the units of its demand in the order $\pi_t$ up to one maximum window, equal for every vessel, and the windows follow one another in the order $(i - t) \bmod M$ with the duration of their grants, with the capacities $\eta(M+1)$ and $\eta(2)$ and without a measurement of the channel.

\begin{table}[t]
\caption{Experiments. Each cell runs 12 seeds; each run lasts 1200 steps.}\label{tab:exp}
\centering\footnotesize
\begin{tabular}{@{}lp{6.3cm}r@{}}
\toprule
& Factors & Runs \\
\midrule
E0 & Calibration: Fleet-DT; S1 to S5; $\rho \in \{1, 1.5, 2\}$; $M = 4$; 4 seeds of their own. Sets $\lambda$, $\gamma$, the initial scores per scenario, and $Z$ of B4. & 60 \\
E1 & Measurement: Fleet-DT; S1 to S5; $\rho = 1.5$; $M = 4$. & 60 \\
E2 & Prediction: Fleet-DT with~\eqref{eq:pred}, persistence, exponential smoothing of $\kappa$, and the calibration mean; S1 to S5; $\rho = 1.5$; $M = 4$. & 240 \\
E3 & Allocation: Fleet-DT, the slotted variant, B1 to B5, and TDMA; S1, S2, S4, S5; $\rho \in \{1, 1.5, 2\}$; $M = 4$. & 1152 \\
E4 & Scale and step: Fleet-DT, B1, B3; S2; $\rho = 1.5$; $M = 2$ with $T_{\mathrm{s}} = 125$~ms, and $M = 4$ with $T_{\mathrm{s}} = 62.5$~ms. & 72 \\
\midrule
E5 & Allocation, addendum: B6 on the cells and the seeds of E3. & 144 \\
E6a & Validation of the simulator: the E3 and E5 run files of Fleet-DT, B1, B3, B6, TDMA; S1, S2; $\rho = 1.5$; $M = 4$; simulated. & 120 \\
E6b & Scale, simulated: Fleet-DT, B1, B3, B6, TDMA; S1, S2; $M = 8$; sources of E3 at $\rho = 1.5$ on every vessel. & 120 \\
E6d & Limit, simulated: Fleet-DT, B1; S2; $M \in \{16, 32\}$; sources of E3 at $\rho = 1.5$ on every vessel. & 48 \\
\bottomrule
\end{tabular}
\end{table}

\subsection{Metrics}\label{subsec:metrics}

The failure fraction $\bar{f}$, the coverage of the bound, the mean absolute error of $\hat{\kappa}$, the ratio of accepted to released bytes per edge, the utilization $\nu$, and the age of information follow Definition~\ref{def:fidelity}, with $\kappa(t)$ from the ground truth. The needed-data delivery ratio $D_{\mathcal{N}}$ divides the bytes of the units of $\mathcal{N}_i(t)$ that the GS accepts within their step by the bytes of the units of $\mathcal{N}_i(t)$ that the sources release, summed over the vessels and the steps of a run. The state miss ratio divides the number of the pairs $(t, i)$ in which the GS misses the state of vessel $i$ within step $t$ by the number of the pairs. The deadline miss ratio $\mu$ is the fraction of the released messages that reach the GS after their step. The work $X(t)$ of Lemma~\ref{lem:delivery} comes from the ground truth: the airtime of the exchanges of the data bridges from $w(t)$ until the data bridges hold no byte of step $t$ or of an earlier step.

\subsection{Procedure and Statistics}\label{subsec:proc}

The calibration runs (E0) use seeds of their own and set: $\lambda$, the gain of~\eqref{eq:pred} with the smallest mean absolute error of $\hat{\kappa}$ in a replay of the E0 runs; the initial scores of each scenario, the scores of its E0 runs; and $\gamma$, the smallest value of $\{0.001, 0.005, 0.01, 0.02, 0.05\}$ for which a replay of~\eqref{eq:aci} on the scores of each E0 run, with the scores of the other E0 runs of its scenario as initial scores and the failure model $f(t) = \mathbb{1}[z(t) < c_S\,Z_t(\alpha_t)]$ with $c_S$ fitted per scenario to the observed failures, gives $\bar{f} \leq \beta + 0.02$ in every run and $|\bar{f} - \beta| \leq 0.02$ in every run in which $\alpha_t = 1$ yields $\bar{f} \geq \beta$. The evaluation runs (E1 to E4) use 12 seeds, and each seed fixes the trajectories, the fading, the interference schedule, the release times and sizes of the sources, and the schedule of the operation control software, so the methods of one cell face the same channel and the same data. A script randomizes the order of the runs. Each run discards its first 240 steps. The tagged commit holds H1 to H7, E1 to E4, the seeds, the parameters of Table~\ref{tab:param}, and the analysis scripts before the first evaluation run. A second tagged commit, made during the runs of E1 to E4 and before any result of E3, adds H8 and E5 (the rows below the rules of Tables~\ref{tab:hyp} and~\ref{tab:exp}); the runs of E5 use the run files of the E3 cells of the same seeds and start after the last run of E4. A third tagged commit adds the simulator, H9, and E6 before the first run of E6, and its revision, after a code review found two defects in the simulator and before any E6 result was read, replaces the runs of the first version; before these commits, the metrics of five E3 runs (Fleet-DT and B1 in S2 at $\rho = 1.5$) set the ACK policy of the simulated TCP, and no comparison between Fleet-DT and a baseline was computed. The archive of the tagged commits and of the data of every run carries a DOI.

The run is the unit of analysis. The comparisons of H3 and H5 pair the runs of one seed and use the Wilcoxon signed-rank test with the Holm correction over all comparisons of the hypothesis, and report the Hodges--Lehmann estimate of the paired difference with its 95\% confidence interval [hollander2013nonparametric, holm1979simple]; H8 uses the same test over the 8 comparisons of each metric. With 12 pairs, the smallest two-sided $p$-value is $2^{-11}$, below the Holm threshold of the largest family, the 40 comparisons of one metric of H5. The confidence interval of $\bar{f}$ within a run uses batch means over 40 batches of 24 steps. H2, H4, and H7 hold for every step or every run, so one counterexample rejects them.

\section{Results}\label{sec:results}
RESULTS (after the campaign). \newpage \mbox{} \newpage \mbox{}
\section{Discussion}\label{sec:discussion}
DISCUSSION. \newpage
\input{back}
\bibliographystyle{IEEEtran}\bibliography{refs_v2,refs_v1sel}
\end{document}
